\section{Faisceaux de morphismes}\etiquette[10]{IV.sec10}\pageoriginale

\begin{enonce*}{Proposition 10.1}\etiquette[10.1]{IV.prop10.1}
Soient $E$ un topos, $X$ et $Y$ deux objets de $E$; le foncteur
$Z\mapsto \Hom_E(Z\times X,Y)\simeq \Hom_{E_{/Z}}(X_Z, Y_Z)$ est
représentable.
\end{enonce*}

En effet, les limites inductives sont universelles dans $E$ (II
4.3). Le foncteur $Z\mapsto Z\times X$ commute donc aux limites
inductives. Par suite, le foncteur $Z\mapsto \Hom_E(Z\times X,Y)$
transforme les limites inductives de l'argument $Z$ en limites
projectives. Il est donc représentable \sisi{(IV \hyperref[IV.prop1.4]{1.4} et
\hyperref[IV.theo1.2]{1.2})}{(\hyperref[IV.cor1.2.1]{1.2.1})}.

\setcounter{subsection}{1}
\subsection{}\etiquette[10.2]{IV.subsec10.2}
L'objet représentant le foncteur $Z\mapsto \Hom_{E}(Z\times X, Y)$
est noté $\scrHom_{E}(X,Y)$ (ou plus simplement $\scrHom(X,Y)$) et
est appelé le \emph{faisceau des morphismes de $X$ dans $Y$}. C'est
un bifoncteur en $X$ et $Y$. On a donc un isomorphisme
trifonctoriel.
\begin{equation*}
\Hom_E(Z,\scrHom_E(X,Y))\simeq \scrHom_E(Z\times X,
Y){\quoi}.\tag{10.2.1}
\end{equation*}\etiquette[10.2.1]{IV.eq10.2.1}

Il résulte alors de la formule \eqhyperref[IV.eq10.2.1]{10.2.1} que le bifoncteur
$(X,Y)\mapsto \scrHom(X,Y)$ transforme les limites inductives de
l'argument $X$ (\resp les limites projectives de l'argument $Y$) en
limites projectives dans $E$.

\begin{enonce*}{Proposition 10.3}\etiquette[10.3]{IV.prop10.3}
Soit $v:E\rightarrow E'$ un morphisme de topos. Pour un objet
variable $X$ de $E$ et un objet variable $Y$ de $E'$, on a un
isomorphisme bifonctoriel.
\begin{equation*}
v_{\ast}\scrHom_E(v^{\ast}Y,X)\simeq
\scrHom_{E'}(Y,v_{\ast}X){\quoi}.\tag{10.3.1}
\end{equation*}\etiquette[10.3.1]{IV.eq10.3.1}
\end{enonce*}

\noindent \emph{Démonstration}. Pour tout objet $Z$ de $E'$, on a
une suite d'isomorphismes, fonctoriels en tout les arguments:
\begin{alignat*}{4}
&\Hom_{E'}(Z,v_{\ast}\scrHom_E(v^{\ast}Y,X)) &&\simeq
&& \Hom_E(v^{\ast}Z,\scrHom_E(v^{\ast}Y,X)) &&\quad \text{(adjonction)}\\
& \Hom_E(v^{\ast}Z,\scrHom_E(v^{\ast}Y,X)) && \simeq
&& \Hom_E(v^{\ast}Z\times v^{\ast}Y,X) &&\quad \eqhyperref[IV.eq10.2.1]{10.2.1}\\
&\Hom_E(v^{\ast}Z\times v^{\ast}Y,X) &&\simeq &&
\Hom_E(v^{\ast}(Z\times Y),X) &&\quad \text{($v^{\ast}$ exact à
gauche)}\\
&\Hom_E(v^{\ast}(Z\times Y),X) &&\simeq && \Hom_{E'}(Z\times Y,
v_{\ast}X) &&\quad \text{(adjonction)}\\
& \Hom_{E'}(Z\times Y, v_{\ast}X) &&\simeq &&
\Hom_{E'}(Z,\scrHom_{E'}(Y,v_{\ast}X)) &&\quad \eqhyperref[IV.eq10.2.1]{10.2.1}
\end{alignat*}
Les deux membres de \eqhyperref[IV.eq10.3.1]{10.3.1} représentent donc des foncteurs
isomorphes, \cqfd

\begin{enonce*}{Corollaire 10.4}\etiquette[10.4]{IV.cor10.4}
Soient\pageoriginale $C$ un $𝒰$-site, $X$ un
$𝒰$-préfaisceau sur $C$, $Y$ un $𝒰$-faisceau sur
$C$, $𝒱$ un univers contenant $𝒰$ tel que $C$ soit
$𝒱$-petit, $C^{\widehat{\phantom{,}}}_{𝒱}$ le topos
des $𝒱$-préfaisceaux sur $C$, $C^{\sim}_{𝒰}$ le
topos des $𝒰$-faisceaux sur $C$. Le préfaisceau
$\scrHom_{C^{\widehat{\phantom{,}}}_{𝒱}}(X,Y)$ est un
$𝒰$-faisceau. Soit $X\rightarrow \underline{a}X$ le morphisme
canonique de $X$ dans son faisceau associé. Le morphisme
correspondant
$\scrHom_{C^{\widehat{\phantom{,}}}_{𝒱}}(\underline{a}X,Y)\rightarrow
\scrHom_{C^{\widehat{\phantom{,}}}_{𝒱}}(X,Y)$ est un
isomorphisme. On a un isomorphisme canonique
$\scrHom_{C^{\widehat{\phantom{,}}}_{𝒱}}(\underline{a}X,Y)\simeq
\scrHom_{C^{\sim}_{𝒰}}(\underline{a}X,Y)$.
\end{enonce*}

\setcounter{subsection}{4} \setcounter{subsubsection}{0}
\subsubsection{}\etiquette[10.4.1]{IV.subsubsec10.4.1}
Supposons d'abord que $C$ soit un petit site et que
$𝒱=𝒰$. On a alors un morphisme de topos:
$C^{\sim}_{𝒰}\rightarrow
C^{\widehat{\phantom{,}}}_{𝒰}$ (\sisi{IV }{}\hyperref[IV.subsec4.8]{4.8}), d'où
\eqhyperref[IV.eq10.3.1]{10.3.1} les assertions dans ce cas. Pour passer de là au cas
général, on remarque tout d'abord que le foncteur « faisceau
associé »  ne dépend pas de l'univers (\hyperrefext[II][II.prop3.6]{3.6}) et que le
topos des $𝒱$-faisceaux sur $C$ est équivalent au topos des
$𝒱$-faisceaux sur $C^{\sim}_{𝒰}$ (\hyperrefext[III][III.sec4]{4} et
\hyperref[IV.sec1]{1}). Il suffit donc de démontrer le lemme suivant:

\begin{enonce*}{Lemme 10.4.2}\etiquette[10.4.2]{IV.lem10.4.2}
Soient $E$ un $𝒰$-topos, $𝒱$ un univers contenant
$𝒰$, $E^{\sim}_{𝒱}$ le topos des
$𝒱$-faisceaux sur $E$, $X$ et $Y$ deux objets de $E$. Il
existe un isomorphisme canonique $\scrHom_E(X,Y)\simeq
\scrHom_{E^{\sim}_{V}}(X,Y)$.
\end{enonce*}

\setcounter{subsubsection}{2}
\subsubsection{}\etiquette[10.4.3]{IV.subsubsec10.4.3}
Le foncteur $\epsilon_{E}:E\rightarrow E_{\widetilde{𝒱}}$
est pleinement fidèle et exact à gauche. Par suite on a, pour tout
objet $Z$ de $E$ un isomorphisme
$$
\Hom_{E_{\widetilde{𝒱}}}(\epsilon_E Z, \epsilon_E
\scrHom_E(X,Y))\simeq \Hom_{E_{\widetilde{𝒱}}}(\epsilon_E
Z\times \epsilon_E X, \epsilon_E Y)
$$
Mais tout objet de $E_{\widetilde{𝒱}}$ est limite inductive
d'objets provenant de $E$, d'où \eqhyperref[IV.eq10.2.1]{10.2.1} l'assertion.

\subsection{}\etiquette[10.5]{IV.subsec10.5}
Soit $v:E\rightarrow E'$ un morphisme de topos. On se propose de
définir quatre morphismes bifonctoriels.
\begin{alignat*}{4}
\Phi_v & : v^{\ast}\scrHom(X,Y) &&\longrightarrow
\scrHom(v^{\ast}X,v^{\ast}Y) &\quav
& X \text{ et } Y \text{ objets de } E'{\quoi},\\
\Psi_v & : v_{\ast}\scrHom(X,Y) &&\longrightarrow \scrHom(v_{\ast}X,
v_{\ast}Y) &\quav & X \text{ et } Y \text{ objets de } E{\quoi},\\
\Xi_v & : \scrHom(v_{!}X,Y) &&\longrightarrow
v_{\ast}\scrHom(X,v^{\ast}Y) &\quav & X\in \ob E, Y\in \ob
E'{\quoi},\\
\Lambda_v & : v_{!}(X\times v^{\ast}Y) &&\longrightarrow
v_{!}(X)\times Y & & X\in \ob E, Y\in \ob E'{\quoi},
\end{alignat*}
où,\pageoriginale pour définir les morphismes $\Xi_v$ et $\Lambda_v$
on suppose que le foncteur $v^{\ast}$ image réciproque par $v$ admet
un adjoint à gauche $v_{!}$.

\subsubsection{}\etiquette[10.5.1]{IV.subsubsec10.5.1}
\emph{Définition de} $\Psi_v$. Soit $X$ un objet de $E$. On a un
morphisme d'adjonction $v^{\ast}v_{\ast}X\rightarrow X$, d'où, pour
tout objet $Z$ de $E'$, un morphisme bifonctoriel $v^{\ast}(Z\times
v_{\ast}X)\rightarrow (v^{\ast}Z)\times X$. On a donc, pour tout
objet $Y$ de $E$, un morphisme trifonctoriel
$$
\Hom_E((v^{\ast}Z)\times X, Y)\longrightarrow
\Hom_E(v^{\ast}(Z\times v_{\ast}X),Y){\quoi}.
$$
On en déduit, par adjonction, un morphisme trifonctoriel
$$
\Hom_{E'}(Z,v_{\ast}\scrHom(X,Y))\longrightarrow \Hom_{E'}(Z,
\scrHom(v_{\ast}X, v_{\ast}Y)){\quoi};
$$
d'où le morphisme $\Psi_v$.

\subsubsection{}\etiquette[10.5.2]{IV.subsubsec10.5.2}
\emph{Définition de} $\Lambda_v$. Pour tout objet $X$ de $E$, on a
un morphisme d'adjonction $X\rightarrow v^{\ast}v_{!}X$, d'où, pour
tout objet $Y$ de $E$, un morphisme bifonctoriel $X\times
v^{\ast}Y\rightarrow v^{\ast}v_{!}(X)\times v^{\ast}Y\simeq
v^{\ast}(v_{!}(X)\times Y)$; d'où, par adjonction, le morphisme
$\Lambda_v$.

\subsubsection{}\etiquette[10.5.3]{IV.subsubsec10.5.3}
\emph{Définition de} $\Xi_v$. Soient $X$ un objet de $E$, $Z$ et $Y$
deux objets de $E'$. Le morphisme $\Lambda_v:v_{!}(X\times
v^{\ast}Z)\rightarrow v_{!}(X)\times Z$ fournit un morphisme
trifonctoriel $\Hom_{E'}(v_{!}(X)\times Z,Y)\rightarrow
\Hom_{E'}(v_{!}(X\times v^{\ast}Z),Y))$; d'où, par adjonction, un
morphisme trifonctoriel
$$
\Hom_{E'}(Z,\scrHom(v_{!}X,Y))\longrightarrow
\Hom_{E'}(Z,v_{\ast}\scrHom(X,v^{\ast}Y)){\quoi};
$$
d'où le morphisme $\Xi_v$.

\subsubsection{}\etiquette[10.5.4]{IV.subsubsec10.5.4}
\emph{Définition de} $\Phi_v$. Le morphisme trifonctoriel de
(\hyperref[IV.subsubsec10.5.3]{10.5.3})
$$
\Hom_{E'}(v_{!}(Z)\times X,Y)\longrightarrow \Hom_{E'}(v_{!}(Z\times
v^{\ast}X),Y){\quoi};
$$
où $Z$ est un objet de $E$ et $X$, $Y$ sont des objets de $E'$,
permet d'obtenir, par adjonction, un morphisme trifonctoriel,
$$
\Hom_E(Z,v^{\ast}\scrHom(X,Y))\longrightarrow
\Hom_E(Z,\scrHom(v^{\ast}X, v^{\ast}Y)){\quoi};
$$
d'où une définition du morphisme $\Phi_v$. Il existe une deuxième
manière de définir $\Phi_v$. Soient $X$, $Y$, $Z$ trois objets de
$E'$. Le foncteur $v^{\ast}:E'\rightarrow E$ fournit un morphisme
trifonctoriel
$$
\Hom_{E'}(Z\times X,Y)\longrightarrow \Hom_E (v^{\ast}(Z)\times
v^{\ast}(X), v^{\ast}Y){\quoi};
$$
d'où, par adjonction, un morphisme trifonctoriel
$$
\Hom_{E'}(Z,\scrHom(X,Y))\longrightarrow
\Hom_{E'}(Z,v_{\ast}\scrHom(v^{\ast}X, v^{\ast}Y)){\quoi}.
$$
On\pageoriginale a donc un morphisme bifonctoriel
$$
\scrHom(X,Y)\longrightarrow v_{\ast}\scrHom(v^{\ast}X,
v^{\ast}Y){\quoi};
$$
d'où, par adjonction, un morphisme dont on laisse au lecteur le soin
de vérifier que c'est bien le morphisme $\Phi_v$ défini ci-dessus.

\begin{enonce*}{Proposition 10.6}\etiquette[10.6]{IV.prop10.6}
Soit $v:E\rightarrow E'$ un morphisme de topos.
\begin{enumerate}
\renewcommand{\labelenumi}{\theenumi)}
\item Le morphisme $\Psi_v$ $(\hyperref[IV.subsubsec10.5.1]{10.5.1})$ est un
  isomorphisme pour tout
objet $Y$ de $E$ si et seulement si le morphisme d'adjonction
$v^{\ast}v_{\ast}X\rightarrow X$ est un isomorphisme. En
particulier, le morphisme $\Psi_v$ est un isomorphisme pour tout
couple d'objets $(X,Y)$ si et seulement si le foncteur
$v_{\ast}:E\rightarrow E'$ est pleinement fidèle, \ie si $v$ est un
plongement de topos.

\item Les conditions suivantes sont équivalentes:
\begin{enumerate}
\renewcommand{\theenumii}{\roman{enumii}}
\renewcommand{\labelenumii}{(\theenumii)}
\item Pour tout couple $(X,Y)$ d'objets de $E'$, le morphisme $\Phi_v$
$(\hyperref[IV.subsubsec10.5.4]{10.5.4})$ est un isomorphisme.

\item Le foncteur $v^{\ast}$ admet un adjoint à gauche $v_{!}$ et pour
tout objet $X$ de $E$ et tout objet $Y$ de $E'$, le morphisme
$\Xi_v$ $(\hyperref[IV.subsubsec10.5.3]{10.5.3})$ est un isomorphisme.

\item Le foncteur $v^{\ast}$ admet un adjoint à gauche $v_{!}$ et pour
tout objet $X$ de $E$ et pour tout objet $Y$ de $E'$, le morphisme
$\Lambda_v$ $(\hyperref[IV.subsubsec10.5.2]{10.5.2})$ est un isomorphisme.
\end{enumerate}
\end{enumerate}
\end{enonce*}

\setcounter{subsection}{6} \setcounter{subsubsection}{0}
\subsubsection{}\etiquette[10.6.1]{IV.subsubsec10.6.1}
La première assertion résulte immédiatement de
(\hyperref[IV.subsubsec10.5.1]{10.5.1}). Il résulte aussi immédiatement de
(\hyperref[IV.subsubsec10.5.2]{10.5.2}), (\hyperref[IV.subsubsec10.5.3]{10.5.3}),
(\hyperref[IV.subsubsec10.5.4]{10.5.4}) que (iii) $\Leftrightarrow$ (ii)
$\Leftrightarrow$ ((i) et le foncteur $v^{\ast}$ admet un adjoint à
gauche $v_{!}$). Il reste donc à montrer que (i) implique que
$v^{\ast}$ admet un adjoint à gauche, \ie (\hyperref[IV.rem1.8]{1.8}) que (i)
implique que $v^{\ast}$ commute aux petits produits. Soient $e'$
l'objet final de $E'$ et $I$ un petit ensemble. Comme $v^{\ast}$ est
exact à gauche, $v^{\ast}(e')$ est un objet final de $E$ et comme
$v^{\ast}$ commute aux limites inductives $v^{\ast}(\amalg_I
e')\simeq \amalg_I v^{\ast}(e')$. Pour tout objet $Y$ de $E'$, on a
$\scrHom(\amalg_I e', Y)\simeq \Pi_I\scrHom(e',Y)\simeq \Pi_I Y$ et
$\scrHom(\amalg_I v^{\ast}(e')$, $v^{\ast}Y)\simeq
\Pi_i\scrHom(v^{\ast}(e'), v^{\ast}Y)\simeq \Pi_I
v^{\ast}Y$.\pageoriginale Le morphisme fonctoriel $\Phi_v$ induit
donc un isomorphisme $v^{\ast}\Pi_I Y\simeq \Pi_I v^{\ast}Y$ dont on
vérifie que c'est l'isomorphisme canonique. \cqfd

\begin{enonce*}{Corollaire 10.7}\etiquette[10.7]{IV.cor10.7}
Soient $E$ un topos, $Z$ un objet de $E$, $j_Z:E_{/Z}\rightarrow E$
le morphisme de localisation (\hyperref[IV.subsec5.2]{5.2}), $X$, $Y$
deux objets de $E$.
\begin{enumerate}
\renewcommand{\labelenumi}{\theenumi)}
\item Il existe un isomorphisme bifonctoriel
$$
j^{\ast}_Z\scrHom(X,Y)\simeq \scrHom(j^{\ast}_ZX,
j^{\ast}_ZY){\quoi}.
$$

\item Soit $X'$ un objet de $E_{/Z}$. Il existe un isomorphisme
bifonctoriel
$$
\scrHom_E(j_{Z!}X',Y)\simeq j_{Z\ast}\scrHom_{E/Z}(X',j^{\ast}_Z
Y){\quoi}.
$$

\item Soit $Y'$ un objet de $E/Z$. Lorsque $Z$ est un ouvert de $E$,
il existe un isomorphisme bifonctoriel
$$
j_{Z\ast}\scrHom(X',Y')\simeq \scrHom(j_{Z\ast}X',
j_{Z\ast}Y'){\quoi}.
$$
\end{enumerate}
\end{enonce*}

Les assertions 1) et 2) se démontrent en remarquant que le morphisme
$\Lambda_{j_Z}$ est un isomorphisme. Pour l'assertion 3), il suffit
de remarquer que $j_{Z\ast}$ est pleinement fidèle, car $j_{Z!}$ est
pleinement fidèle (\hyperrefext[I][I.rem5.7]{5.7}. a)).

\begin{enonce*}{Corollaire 10.8}\etiquette[10.8]{IV.cor10.8}
Soient $E$ un topos, $X$, $Y$ et $Z$ trois objets de $E$,
$j_Z:E_{/Z}\rightarrow E$ le morphisme de localisation.
\begin{enumerate}
\renewcommand{\labelenumi}{\theenumi)}
\item On a un isomorphisme canonique
$$
j_{Z\ast}j^{\ast}_ZY\simeq \scrHom(Z,Y){\quoi}.
$$

\item On a un isomorphisme canonique
$$
\Hom_E(Z,\scrHom(X,Y))\simeq \Hom_{E_{/Z}}(j^{\ast}_ZX,
j^{\ast}_ZY){\quoi}.
$$
\end{enumerate}
\end{enonce*}

Soit $e_Z$ l'objet final de $E_{/Z}$ (\ie l'objet
$\id_Z:Z\rightarrow Z$)). Il résulte de \eqhyperref[IV.eq10.2.1]{10.2.1} qu'on a un
isomorphisme canonique
$$
\scrHom_{E/Z}(e_Z, j^{\ast}_ZY)\simeq j^{\ast}_ZY{\quoi};
$$
d'où, d'après (\hyperref[IV.cor10.7]{10.7} 2), un isomorphisme
$$
j_{Z\ast}j^{\ast}_Z Y\simeq \scrHom(j_{Z!}e_Z, Y){\quoi.}
$$
Comme $j_{Z!}e_Z=Z$, on a démontré 1). Démontrons 2). De
(\hyperref[IV.cor10.7]{10.7}), on tire un\pageoriginale isomorphisme canonique
$$
\Hom_{E_{/Z}}(e_Z, j^{\ast}_Z\scrHom(X,Y))\simeq
\Hom_{E_{/Z}}(j^{\ast}_Z X, j^{\ast}_ZY){\quoi};
$$
d'où, par adjonction sur le premier membre,
$$
\Hom_E(Z,\scrHom(X,Y))\simeq \Hom_{E_{/Z}}(j^{\ast}_ZX,
j^{\ast}_ZY){\quoi}.
$$

\section{Topos annelés, localisation dans les topos annelés}\etiquette[11]{IV.sec11}

\setcounter{subsection}{1} \setcounter{subsubsection}{0}
\subsubsection{}\etiquette[11.1.1]{IV.subsubsec11.1.1}
Soit $𝒰$ un univers. On appelle $𝒰$-\emph{topos
annelé} un couple $(E,A)$ où $E$ est un $𝒰$-topos et $A$ un
objet muni d'une structure d'anneau. On appelle $𝒰$-site
annelé un couple $(C,A)$ où $C$ est un $𝒰$-site et $A$ un
$𝒰$-faisceaux d'anneaux sur $C$. On ne mentionne pas
l'univers lorsque le contexte ne prête pas à confusion. À un site
annelé $(C,A)$ est associé le topos annelé $(C^{\sim},A)$. À un
topos annelé $(E,A)$ est associé le site annelé
constitué par le
site $E$ et le faisceau d'anneaux représenté par $A$.

\subsubsection{}\etiquette[11.1.2]{IV.subsubsec11.1.2}
Soit $(E,A)$ un topos annelé. On note ${}_AE$ (\resp $E_A$) la
catégorie des faisceaux de $A$-modules (nous écrirons aussi
$A$-Modules) à gauche (\resp à droite) unitaires. La catégorie
${}_AE$ (\resp $E_A$) est une catégorie abélienne
(\hyperrefext[II][II.prop6.7]{6.7}). Soient $M$ et $N$ deux faisceaux de $A$-modules
à gauche (\resp à droite). Le groupe commutatif des morphismes de
$A$-Modules de $M$ dans $N$ est noté $\Hom_A(M,N)$.

\subsubsection{}\etiquette[11.1.3]{IV.subsubsec11.1.3}
Soient $A$, $B$ et $C$ trois anneaux d'un topos $E$, $M$ un faisceau
de $A-B$ bimodules, $N$ un faisceau de $A-C$ bimodules à gauche.
Soit $e$ un objet final de $E$ et posons $\Hom(e,B)=\Gamma(B)$,
$\Hom(e,C)=\Gamma(C)$. La structure de $A-B$ bimodule de $M$ fournit
un homomorphisme d'anneaux de $\Gamma(B)^{\circ}$ dans l'anneau des
endomorphismes du $A$-module $M$ et de même, la structure de $A-C$
bimodule de $N$ fournit un homomorphisme d'anneaux de
$\Gamma(C)^{\circ}$ dans l'anneau des endomorphismes du $A$-module
$N$. On en déduit, par fonctorialité, une structure de
$\Gamma(B)-\Gamma(C)$ bimodule sur le groupe commutatif
$\Hom_A(M,N)$. En particulier, lorsque $A$ est un faisceau d'anneaux
commutatifs, le groupe $\Hom_A(M,N)$ est muni canoniquement d'une
structure de $\Gamma(A)$-module.

\subsubsection{}\etiquette[11.1.4]{IV.subsubsec11.1.4}
Soient\pageoriginale $E$ un topos, $A$ et $B$ deux anneaux de $E$,
$u:A\rightarrow B$ un morphisme de faisceaux d'anneaux, $M$ un
$B$-module (à gauche pour fixer les idées). Le faisceau $M$
peut-être considéré comme un faisceau de $A$-modules par
l'intermédiaire de $u$: pour tout objet $X$ de $E$, $M(X)$ est muni
de la structure de $A(X)$-module déduite de sa structure de
$B(X)$-module et de l'homomorphisme $u(X):A(X)\rightarrow B(X)$ par
restriction des scalaires. On obtient ainsi un foncteur « \emph{restriction des scalaires par} $u$ »
$$
\Res(u):{}_BE\longrightarrow {}_AE{\quoi}.
$$

\subsubsection{}\etiquette[11.1.5]{IV.subsubsec11.1.5}
En particulier lorsque $A=\ZZ$ (faisceau constant $\sisi{\underline{\underline{\ZZ}}}{\underline{\ZZ}}$ \ie faisceau
associé au préfaisceau constant $\ZZ$) et lorsque $u:\sisi{\underline{\underline{\ZZ}}}{\underline{\ZZ}}\rightarrow B$
est l'unique morphisme canonique, on obtient un foncteur de ${}_BF$
dans la catégorie ${}_{\sisi{\underline{\underline{\ZZ}}}{\underline{\ZZ}}}E$ qui n'est autre que la catégorie des
faisceaux abéliens notée $E_{Ab}$. Ce foncteur est appelé le
\emph{foncteur faisceau abélien sous-jacent}.

\begin{enonce*}{Proposition 11.1.6}\etiquette[11.1.6]{IV.prop11.1.6}
Le foncteur restriction des scalaires par $u$ commute aux limites
inductives et projectives. Il est conservatif.
\end{enonce*}

La proposition est vraie pour le foncteur restriction des scalaires
pour les modules ordinaires, \ie lorsque $E$ est le topos ponctuel.
Elle est donc vraie lorsque $E$ est le topos des préfaisceaux sur un
petit site. Il résulte alors de la détermination des limites
inductives et projectives à l'aide du foncteur faisceau associé
(\hyperrefext[II][II.prop6.4]{6.4}) que la proposition est vraie dans le cas général.

\setcounter{subsection}{2} \setcounter{subsubsection}{0}
\subsubsection{}\etiquette[11.2.1]{IV.subsubsec11.2.1}
Soient $(E,A)$ un topos annelé, $X$ un objet de $E$,
$j_X:E_{/X}\rightarrow E$ le morphisme de localisation
\sisi{(\hyperref[IV.sec8]{8})}{\hyperref[IV.subsec5.2]{5.2}}. D'après \hyperrefext[III][III.prop1.7]{1.7} ou
\hyperref[IV.subsubsec3.1.2]{3.1.2}, le faisceau $j^{\ast}_XA$ est muni
canoniquement d'une structure d'anneau. Le faisceau d'anneaux
$j^{\ast}_XA$ est noté le plus souvent $A| X$ ou bien encore,
abusivement, $A$. Sauf mention du contraire, le topos $E_{/X}$ sera
annelé par $A| X$. Le faisceau $j_{X\ast}j^{\ast}_XA$ est muni
canoniquement d'une structure d'anneau. Le morphisme d'adjonction
$A\rightarrow j_{X\ast}j^{\ast}_XA$ est un morphisme de faisceaux
d'anneaux.

\subsubsection{}\etiquette[11.2.2]{IV.subsubsec11.2.2}
Soit\pageoriginale de plus $M$ un $A$-module (à gauche pour fixer
les idées). Le faisceau $j^{\ast}_XM$ est muni d'une structure de
$A| X$-module (\hyperrefext[III][III.prop1.7]{1.7}; ou \hyperref[IV.subsubsec3.1.2]{3.1.2}); d'où
un foncteur:
$$
j^{\ast}_X:{}_AE\longrightarrow {}_{A| X}E_{/X}
$$
appelé foncteur de restriction à $E_{/X}$, ou encore foncteur de
restriction à $X$. Le foncteur $j^{\ast}_X$ commute aux limites
inductives et projectives (\loccit). Il est en particulier exact.
Soit maintenant $N$ un $A| X$-Module. Le faisceau $j_{X\ast}N$ est
un faisceau de $j_{X\ast}A| X$-Modules (\hyperrefext[III][III.prop1.7]{1.7} ou 
\hyperref[IV.subsubsec3.1.2]{3.1.2}); d'où, par restriction des scalaires par le
morphisme d'adjonction $A\rightarrow j_{X\ast}(A| X)$
(\hyperref[IV.subsubsec11.1.4]{11.1.4}), un $A$-Module encore noté
$j_{X\ast}N$. On a donc défini un foncteur
$$
j_{X\ast}:{}_{A| X}E_{/X}\longrightarrow {}_AE
$$
qui commute aux limites projectives (\hyperref[IV.prop11.1.6]{11.1.6} et
\hyperrefext[III][III.prop1.7]{1.7}). Pour tout $A$-Module $M$, le morphisme
d'adjonction $M\rightarrow j_{X\ast}j^{\ast}_XM$ est un morphisme de
$A$-Modules. Pour tout $A| X$-Module $N$ et tout $A$-Module $M$, le
morphisme d'adjonction $M\rightarrow j_{X\ast}j^{\ast}_XM$ définit
un morphisme bifonctoriel en $M$ et $N$
\begin{equation*}
\Hom_{A| X}(j^{\ast}_XM,N)\longrightarrow
\Hom_A(M,j_{X\ast}N){\quoi}.\tag{11.2.2.1}
\end{equation*}\etiquette[11.2.2.1]{IV.eq11.2.2.1}
Ce dernier morphisme est un isomorphisme \ie les foncteurs
$j^{\ast}_X$ et $j_{X\ast}$ pour les $A$-Modules, sont adjoints (III
1.7).

\subsubsection{}\etiquette[11.2.3]{IV.subsubsec11.2.3}
Les foncteurs $j^{\ast}_X$ et $j_{X\ast}$ pour les Modules,
commutent au foncteur « ensemble sous-jacent »
(\hyperrefext[III][III.prop1.7]{1.7} 4)). Ils commutent donc aux foncteurs restriction
des scalaires et en particulier au foncteur faisceau abélien
sous-jacent.

\setcounter{subsection}{3} \setcounter{subsubsection}{0}
\begin{enonce*}{Proposition 11.3.1}\etiquette[11.3.1]{IV.prop11.3.1}
Soient $(E,A)$ un topos annelé, $X$ un objet de $E$. Le foncteur
$$
j^{\ast}_X:{}_{A}E\longrightarrow {}_{A| X}E_{/X}
$$
admet un adjoint à gauche noté $j_{X!}:{}_{A| X}E_{/X}\rightarrow
{}_{A}E$ et appelé le prolongement par zéro. Le foncteur
prolongement par zéro est exact et fidèle et commute aux limites
inductives. Les foncteurs $j_{X!}$ pour les Modules\pageoriginale
commutent aux foncteurs restriction des scalaires et, en
particulier, ils commutent au foncteur faisceau abélien sous-jacent.
\end{enonce*}

L'existence du foncteur $j_{X!}$ résulte de \hyperrefext[III][III.prop1.7]{1.7}. Comme
$j_{X!}$ est un adjoint à gauche, il commute aux limites inductives
(\hyperrefext[I][I.sec2]{2}). Pour démontrer les autres assertions, supposons
d'abord que $E$ soit le topos des préfaisceaux d'ensembles sur une
petite catégorie $C$ contenant l'objet $X$. On sait alors que
$E_{/X}$ est équivalent à la catégorie des préfaisceaux sur $C_{/X}$
et que, modulo cette équivalence, le foncteur
$j^{\ast}_X:E\rightarrow E_{/X}$ n'est autre que la composition avec
le foncteur d'oubli $C_{/X}\rightarrow C$ (\hyperrefext[I][I.prop5.11]{5.11}). Il
résulte alors immédiatement de la construction explicite de $j_{X!}$
(\hyperrefext[I][I.prop5.1]{5.1}) que pour tout $A| X$-Module $N$ et pour tout objet
$Y$ de $C$, on a
$$
j_{X!}N(Y)=\bigoplus_{u\in \Hom_C(Y,X)}N(u){\quoi};
$$
d'où l'exactitude de $j_{X!}$ et le fait que le prolongement par
zéro commute aux foncteurs restriction des scalaires dans ce cas.
Dans le cas général, on peut supposer que $E$ est le topos des
faisceaux sur un petit site $C$ et que $X$ provient d'un objet de
$C$ (\hyperref[IV.cor1.2.1]{1.2.1}). Le topos $E_{/X}$ est alors équivalent au
topos des faisceaux sur $C_{/X}$ (muni de la topologie induite)
(\hyperrefext[III][III.prop5.4]{5.4}) et le morphisme de topos $j_X:E_{/X}\rightarrow
E$ provient du foncteur d'oubli $C_{/X}\rightarrow C$ qui est
continu et cocontinu (\hyperrefext[III][III.prop5.2]{5.2}). Il résulte alors de
\hyperrefext[III][III.prop1.7]{1.7}) que le prolongement par zéro pour les
faisceaux s'obtient en composant le prolongement par zéro pour les
préfaisceaux avec le foncteur faisceau associé; d'où l'assertion
d'exactitude et la commutation aux restrictions des scalaires.

Pour démontrer la fidélité, il revient au même de
montrer que pour
$A| X$-Module $N$, le morphisme d'adjonction
$$
\ad_N:N\longrightarrow j^{\ast}_X j_{X!}N
$$
est un monomorphisme. Or, en notant $j_{\widehat{X}!}$ le foncteur
prolongement par zéro pour les préfaisceaux et $\underline{a}$ le
foncteur « faisceau associé », on a un diagramme commutatif
\[
\xymatrix@C=2.75cm@R=1.5cm{ N \ar[r]^{\ad_{N}}
\ar[rd]_{\ad_{\hat{N}}} & j^{\ast}_{X} \underline{a}
j_{\widehat{X}!} N \\
                   & \underline{a} j^{\ast}_X j_{\hat{X!}} N,\ar[u]
}
\]
où\pageoriginale $\ad_{\widehat{N}}$ est le morphisme d'adjonction
pour les préfaisceaux. Le morphisme $\ad_{\widehat{N}}$ est un
monomorphisme d'après ce qui précède, donc
$\underline{a}(\ad_{\widehat{N}})$ est un monomorphisme. De plus, comme
le foncteur d'oubli $C_{/X}\rightarrow C$ est continu et cocontinu,
le morphisme canonique $\underline{a}j^{\ast}_X\rightarrow
j^{\ast}_X\underline{a}$ est un isomorphisme (\hyperrefext[III][III.prop2.3]{2.3}). Par
suite $\ad_N$ est un monomorphisme.

\begin{enonce*}[remark]{Remarque 11.3.2}\etiquette[11.3.2]{IV.rem11.3.2}
Soit $N$ un $A| X$-Module. Le faisceau d'ensemble sous-jacent à
$j_{X!}N$ n'est pas, en général, isomorphe au faisceau obtenu en
prolongeant par le vide (\ref{III.subsec5.3} et \hyperref[IV.subsec5.2]{5.2}) le faisceau
d'ensemble sous-jacent à $N$. On prendra donc garde de ne pas
confondre le foncteur \emph{prolongement par zéro} noté
$j_{X!}:{}_{A|X}E_{/X}\rightarrow {}_{A}E$ dans \hyperref[IV.prop11.3.1]{11.3.1},
et le foncteur prolongement par le vide noté encore
$j_{X!}:E_{/X}\rightarrow E$ dans III 5.3 et IV 5.2. Dans la plupart
des cas rencontrés dans la pratique, l'abus de notations signalé
ci-dessus n'amène pas de confusions. Lorsqu'une confusion est
néanmoins possible, nous utiliserons les notations $j^{\ab}_{X!}$ et
$j^{\ens}_{X!}$ pour désigner respectivement les foncteurs
prolongement par zéro et prolongement par le vide.
\end{enonce*}

\begin{enonce*}{Proposition 11.3.3}\etiquette[11.3.3]{IV.prop11.3.3}
Soient $(E,A)$ un topos annelé et $X$ un objet de $E$. Le $A$-Module
$j_{X!}(A| X)$, noté le plus souvent $A_X$ ou $A_{X,E}$, est le
$A$-Module libre engendré par $X$ {\rm (\hyperrefext[II][II.prop6.5]{6.5})} \ie pour
tout $A$-Module $M$, on a un isomorphisme canonique, fonctoriel en
$M$:
$$
\Hom_E(X,M)\xleftarrow{\sim} \Hom_A (A_X, M){\quoi}.
$$
Soit $(X_i)_{i\in I}$ une famille topologiquement génératrice de $E$
{\rm (\hyperrefext[II][II.def3.0.1]{3.0.1}).} La famille $(A_{X_i})_{i\in I}$ est une
famille génératrice de la catégorie des $A$-Modules.
\end{enonce*}

Soit $e_x$ l'objet final du topos $E_{/X}$
$(e_X=X\xrightarrow{\id}X)$. On a un isomorphisme
$$
\Hom_{A| X}(A| X, j^{\ast}_X M)\xrightarrow{\sim} \Hom_{E_{/X}}(e_X,
j^{\ast}_X M){\quoi};
$$
déduit\pageoriginale de la section unité $e_X\rightarrow A/X$; d'où,
par adjonction, un isomorphisme
$$
\Hom_A(j^{\ab}_{X!}(A| X),
M)\xrightarrow{\sim}\Hom_E(j^{\ens}_{X!}(e_X),M){\quoi}.
$$
Comme $j^{\ens}_{X!}(e_X)=X$, on obtient l'isomorphisme annoncé. La
dernière assertion résulte de \hyperrefext[II][II.cor6.6]{6.6}.

\begin{enonce*}[remark]{Remarque 11.3.4}\etiquette[11.3.4]{IV.rem11.3.4}
Soient $A$ et $B$ deux faisceaux d'anneaux sur un topos $E$, $X$ un
objet de $E$ et $N$ un $A|X$-$B| X$-bi Module. Le faisceau abélien
obtenu en prolongeant $N$ par zéro est muni canoniquement d'une
structure de $A$-$B$-biModule ainsi qu'il résulte immédiatement de
sa description explicite (\hyperref[IV.prop11.3.1]{11.3.1}). En particulier le
$A$-module libre engendré $A_X$ est un $A$-biModule.
\end{enonce*}

\begin{enonce*}[remark]{Exercice 11.3.5}\etiquette[11.3.5]{IV.exer11.3.5}
Soient $E$ un topos, $X$ un objet de $E$, $x:P\rightarrow E$ un
point de $E$ (\hyperref[IV.def6.1]{6.1}), $(x_i)_{i\in I}$ la famille des
points de $E_{/X}$ au-dessus de $x$ (en correspondance biunivoque
avec la fibre $X_x$ (\hyperref[IV.eq6.7.2]{6.7.2}). Montrer que pour tout
faisceau abélien $M$ sur $E_{/X}$, la fibre $(j_{X!}M)_x$ est
canoniquement isomorphe à $\oplus_i M_{x_i}$.
\end{enonce*}

\section{Opération sur les modules}\etiquette[12]{IV.sec12}

\begin{enonce*}{Proposition 12.1}\etiquette[12.1]{IV.prop12.1}
Soient $(E,A)$ un topos annelé, $M$ et $N$ deux $A$-Modules à gauche
(\resp à droite). Le foncteur sur $E$ qui a tout objet $X$ de $E$
associe le groupe commutatif $\Hom_A(M,\scrHom_E(X,N))$ est
représentable par un faisceau abélien noté $\scrHom_A(M,N)$ (ou
parfois $\scrHom(M,N)$ lorsqu'aucune confusion n'en résulte). Pour
tout objet $X$ de $E$ on a un isomorphisme canonique
\begin{equation*}
\scrHom_A(M,N)(X)\simeq \Hom_{A| X}(j^{\ast}_XM,
j^{\ast}_XN){\quoi}.\tag{12.1.1}
\end{equation*}\etiquette[12.1.1]{IV.eq12.1.1}
\end{enonce*}

On a un isomorphisme canonique $\scrHom_E(X,N)=j_{X\ast}j^{\ast}_XN$
(\hyperref[IV.cor10.8]{10.8}) et par suite $\scrHom_E(X,N)$ est muni
fonctoriellement en $X$ d'une structure de $A$-Module à gauche
(\resp à droite). Le foncteur $X\rightarrow \scrHom_E(X,N)$
transforme les limites inductives de l'argument $X$ en limites
projectives de $A$-Modules (\hyperref[IV.subsec10.2]{10.2} et
\hyperref[IV.subsubsec11.2.3]{11.2.3}) et par suite le foncteur
$$
X\longrightarrow \Hom_Z(M,\scrHom_E(X,N))
$$
transforme les limites inductives\pageoriginale de l'argument $X$ en
limites projectives de groupes commutatifs, donc en limites
projectives des ensembles sous-jacents. Il est donc représentable
(\hyperref[IV.prop1.4]{1.4} et \hyperref[IV.theo1.2]{1.2}) et l'objet qui le représente
est muni d'une structure de faisceau abélien. On a, par définition,
pour tout objet $X$ de $E$, un isomorphisme canonique
\begin{equation*}
\scrHom_A(M,N)(X)=\Hom(X,\scrHom_A(M,N))\simeq
\Hom_A(M,\scrHom(X,N))\tag{12.1.2}
\end{equation*}\etiquette[12.1.2]{IV.eq12.1.2}
et l'isomorphisme (\hyperref[IV.eq12.1.1]{12.1.1}) résulte de \eqhyperref[IV.eq12.1.2]{12.1.2}, de
l'isomorphisme $\scrHom(X,N)\approx j_{X\ast}j^{\ast}_XN$
(\hyperref[IV.cor10.8]{10.8}) et des formules d'adjonction de 10.

\setcounter{subsection}{1}
\subsection{}\etiquette[12.2]{IV.subsec12.2}
Le faisceau abélienne $\scrHom_A(M,N)$ est appelé le \emph{faisceau
des morphismes de $A$-Modules de $M$ dans $N$}. C'est un bifoncteur
en $M$ et $N$. Il résulte de sa définition et de
(\hyperref[IV.subsec10.2]{10.2}) qu'il transforme les limites projectives de
l'argument $N$ (\resp les limites inductives de $M$) en limites
projectives de faisceaux abéliens. En particulier il est exact à
gauche en ses deux arguments.

\begin{enonce*}{Proposition 12.3}\etiquette[12.3]{IV.prop12.3}
Soient $(E,A)$ un topos annelé, $M$ et $N$ deux $A$-Modules à droite
(\resp à gauche), $X$ un objet de $E$:
\begin{enumerate}
\renewcommand{\theenumi}{\alph{enumi}}
\renewcommand{\labelenumi}{\theenumi)}
\item On a des isomorphismes canoniques:
$$
\scrHom_A(M,N)(X)\approx \Hom_A(M, j_{X\ast}j^{\ast}_XN)\simeq
\Hom_{A| X}(j^{\ast}_XM, j^{\ast}_XN)\simeq
\Hom_A(j_{X!}j^{\ast}_XM, N)
$$

\item On a un isomorphisme canonique:
$$
\Phi_X:j^{\ast}_X\scrHom_A (M,N)\simeq \scrHom_{A| X}(j^{\ast}_X M,
j^{\ast}_XN)
$$
Soit de plus $P$ un $A| X$-Module à droite (\resp à gauche):

\item On a un isomorphisme canonique:
$$
j_{X\ast}\scrHom_{A| X}(j^{\ast}_X M,P)\simeq
\scrHom_A(M,j_{X\ast}P)
$$

\item On a un isomorphisme canonique:
$$
\Xi_X:\scrHom_A (j_{X!}P,N)\simeq j_{X\ast}\scrHom_{A|
X}(P,j^{\ast}_XN)
$$
\end{enumerate}
\end{enonce*}

\noindent Les isomorphismes de a) résultent de \eqhyperref[IV.eq12.1.2]{12.1.2}, de
l'isomorphisme $\scrHom_E(X,N)\simeq j_{X\ast}j^{\ast}_XN$
(\hyperref[IV.cor10.8]{10.8}) et des formules d'adjonction de 10.

\begin{itemize}
\item[Démontrons d).] Pour tout objet $Y$ de $E_{/X}$ on a la suite
d'isomorphismes:

$\Hom(Y,j^{\ast}_X\scrHom_A(M,N))\simeq \Hom(j_{X!}Y,\scrHom_A
(M,N))$ (adjonction pour les faisceaux d'ensembles)
\begin{alignat*}{2}
& \Hom(j_{X!}Y,\scrHom_A(M,N))\simeq \Hom_A (M,\scrHom_E(j_{X!}Y,
N))
&\quad& (12.2.1)\\
& \Hom(M, \scrHom_E(j_{X!}Y,N))\simeq \Hom_A
(M,j_{X\ast}\scrHom_{E_{/X}}(Y,j^{\ast}_XN)) && (\hyperref[IV.cor10.7]{10.7})\\
& \Hom_{\Lambda}(M, j_{X\ast}\scrHom_{E_{/X}}(Y, j^{\ast}_XN))\simeq
\Hom_{A/X}(j^{\ast}_XM, \scrHom_E (Y,j^{\ast}_XN)) && \eqhyperref[IV.eq11.2.2.1]{11.2.2.1}\\
 & \Hom_{A| X}(j^{\ast}_XM, \scrHom_E (Y, j^{\ast}_XN))\simeq
\Hom(Y,\scrHom_{A| X}(j^{\ast}_X M, j^{\ast}_XM)) && (12.2.1)
\end{alignat*}\pageoriginale

\item[Démontrons c).] Pour tout objet $Y$ de $E$, on a la suite
d'isomorphismes:
\begin{alignat*}{2}
& \Hom(Y,j_{X\ast}\scrHom_{A| X}(j^{\ast}_XM,P))\simeq
\Hom(j^{\ast}_XY, \scrHom_{A| X}(j^{\ast}_XM; P)) &\quad&
\text{(adjonction)}\\
&\Hom (j^{\ast}_XY, \scrHom_{A| X}(j^{\ast}_XM, P))\simeq
\Hom_{A/X}(j^{\ast}_XM, \scrHom_{E_{/X}}(j^{\ast}_XY,P)) && (12.2.1)\\
&\Hom_{A| X}(j^{\ast}_XM, \scrHom_{E_{/X}}(j^{\ast}_XY,P))\simeq
\Hom_A(M, j_{X\ast}\scrHom_{E_{/X}}(j^{\ast}_XY,P)) && \eqhyperref[IV.eq11.2.2.1]{11.2.2.1}\\
&\Hom_A(M,j_{X\ast}\scrHom_{E_{/X}}(j^{\ast}_XY,P))\simeq
\Hom_A(M,\scrHom_E(Y, j_{X\ast}P)) && (\hyperref[IV.prop10.3]{10.3})\\
&\Hom_A(M,\scrHom_E(Y, j_{X\ast}P))\simeq \Hom(Y,
\scrHom_A(M,j_{X\ast}P)) && (12.2.1)
\end{alignat*}

\item[Démontrons d).] Pour tout objet $Y$ de $E$, on à la suite
d'isomorphismes:
\begin{alignat*}{2}
& \Hom(Y,\scrHom_A(j_{X!}P,N))\simeq \Hom_A(j_{X!}P,\scrHom_E(Y,N))
&\quad& (12.2.1)\\
&\Hom_A(j_{X!}P,\scrHom_E(Y,N))\simeq \Hom_{A|
X}(P,j^{\ast}_X\scrHom_E(Y,N)) && (\hyperref[IV.prop11.3.1]{11.3.1})\\
&\Hom_{A/X}(P,j^{\ast}_X\scrHom_E(Y,N))\simeq \Hom_{A|
X}(P,\scrHom_{E_{/X}}(j^{\ast}_XY, j^{\ast}_XN)) && (\hyperref[IV.cor10.7]{10.7})\\
&\Hom_{A/X}(P,\scrHom_{E_{/X}}(j^{\ast}_XY, j^{\ast}_XN))\simeq
\Hom(j^{\ast}_XY, \scrHom_{A/X}(P,j^{\ast}_XN)) && (12.2.1)\\
&\Hom(j^{\ast}_XY, \scrHom_{A| X}(P,j^{\ast}_XN))\simeq (Y,
j_{X\ast}\scrHom_{A| X}(P, j^{\ast}_XN)) && \text{(adjonction)}
\end{alignat*}
\end{itemize}

\begin{enonce*}{Corollaire 12.4}\etiquette[12.4]{IV.cor12.4}
Soient $C$ un $𝒰$-site, $A'$ un $𝒰$-préfaisceau
d'anneaux sur $C$, $M$ un $𝒰$-préfaisceau de $A'$-modules (à
gauche pour fixer les idées), $N$ un $𝒰$-faisceau de
$A'$-modules (à gauche). Notons $A$ le faisceau associé à $A'$, de
sorte que les faisceaux $N$ et $\underline{a}M$ (faisceau associé à $M$)
sont des $A$-Modules. Le préfaisceau $X\mapsto \Hom_{A'|
X}(j^{\ast}_XM, j^{\ast}_XN)$ est un $𝒰$-faisceau abélien
canoniquement isomorphe à $\scrHom_A(\underline{a}M,N)$.
\end{enonce*}

En effet il résulte de \hyperrefext[III][III.prop5.5]{5.5} et \hyperrefext[III][III.prop2.3]{2.3} que
le foncteur de localisation à $C_{/X}$ commute avec les foncteurs
faisceaux associés (pour $C$ et $C_{/X}$). Par suite, on a des
isomorphismes fonctoriels en l'objet variable $X$ de $C$:
\begin{align*}
\Hom_{A'| X}(j^{\ast}_XM, j^{\ast}_XN) &\simeq \Hom_{A|
X}(\underline{a}j^{\ast}_XM, j^{\ast}_XN)\\
&\simeq \Hom_{A| X}(j^{\ast}_X \underline{a}M, j^{\ast}_XN)\simeq
\scrHom_A (\underline{a}M, N)(X){\quoi}.
\end{align*}

\begin{enonce*}{Corollaire 12.5}\etiquette[12.5]{IV.cor12.5}
Soient\pageoriginale $E$ un topos, $A$, $B$, $C$ trois faisceaux
d'anneaux sur $E$, $M$ un $A-B$ biModule, $N$ un $A-C$ biModule. Le
faisceau abélien $\scrHom_A(M,N)$ est muni canoniquement d'une
structure de $B-C$ biModule. En particulier lorsque $A$ est un
faisceau d'anneaux commutatifs et lorsque $M$ et $N$ sont des
$A$-Modules, $\scrHom_A(M,N)$ est muni canoniquement d'une structure
de $A$-Module. Les isomorphismes canoniques {\rm b), c), d)} de
$\hyperref[IV.prop12.3]{12.3}$ sont des isomorphismes de biModules.
\end{enonce*}

Nous nous bornerons à donner des indications. Pour tout objet $X$ de
$E$, on a $\scrHom_A(M,N)(X)\simeq \Hom_{A| X}(j^{\ast}_XM,
j^{\ast}_XN)$ (\hyperref[IV.prop12.3]{12.3}). Par suite, le groupe commutatif\\
$\scrHom_A(M,N)(X)$ est muni canoniquement d'une structure de
$B(X)-C(X)$ bimodule (\hyperref[IV.subsubsec11.1.3]{11.1.3}) dont on vérifie
qu'elle est fonctorielle en $X$. Les autres assertions sont laissées
au lecteur.

\begin{enonce*}{Corollaire 12.6}\etiquette[12.6]{IV.cor12.6}
Soient $(E,A)$ un topos annelé, $X$ un objet de $A$, $N$ un
$A$-Module (à gauche pour fixer les idées). On a des isomorphismes
canoniques de $A$-Modules:
$$
\scrHom_E(X,N)\simeq j_{X\ast}j^{\ast}_XN\simeq \scrHom_A(A_X,
N){\quoi};
$$
la structure de $A$-Module sur $\scrHom_A(A_X,N)$ provenant de la
structure de biModule sur $A_X$ $(\hyperref[IV.rem11.3.4]{11.3.4})$.
\end{enonce*}

Le premier isomorphisme résulte de \hyperref[IV.cor10.7]{10.7}, le deuxième de
l'isomorphisme de $A$-Modules $N\simeq \scrHom_A(A,N)$ et de
\hyperref[IV.prop12.3]{12.3}.

\begin{enonce*}{Proposition 12.7}\etiquette[12.7]{IV.prop12.7}
Soient $(E,A)$ un topos annelé, $M$ un $A$-Module à droite et $N$ un
$A$-Module à gauche. Le foncteur qui à tout faisceau abélien $P$
associe le groupe commutatif
$\Hom_A(M,\scrHom_{\sisi{\underline{\underline{\ZZ}}}{\underline{\ZZ}}}(N,P))$ est représentable par
un faisceau abélien noté $M\otimes_A N$ et appelé le produit
tensoriel sur $A$ de $M$ et de $N$.
\end{enonce*}

On a une injection canonique de $\Hom_A(M,\scrHom_{\sisi{\underline{\underline{\ZZ}}}{\underline{\ZZ}}}(N,P))$ dans
l'ensemble \\ $\Hom(M,\scrHom(N,P))\simeq \Hom(M\times N, P)$; On
constate aussitôt, en revenant aux définitions, qu'un morphisme $f$
de faisceaux d'ensembles de $M\times N$ dans $P$ provient d'une
élément de $\Hom_A(M,\scrHom_Z(N,P))$ si et seulement si pour tout
objet $X$ de $E$, $f(X):M(X)\times N(X)\rightarrow P(X)$ est une
application\pageoriginale $A(X)$-bilinéaire de $M(X)\times N(X)$
dans $P(X)$. Appelons morphisme $A$-bilinéaire les morphismes de
faisceaux d'ensembles de $M\times N$ dans $P$ qui possèdent cette
propriété. II s'agit donc de représenter le foncteur des morphismes
$A$-bilinéaires de $M\times N$ dans $P$. Pour cela on procède comme
dans le cas ordinaire \ie comme dans le cas où $E$ est le topos
ponctuel. Considérons les huit morphismes:
\begin{itemize}
\item[(i):] $M\times M\times N \longrightarrow M\times N$\qquad $1\leq
i\leq 3$

\item[(i):] $M\times N\times N\longrightarrow M\times N$\qquad $4\leq
i\leq 6$

\item[(i):] $M\times A\times N\longrightarrow M\times N$\qquad $7\leq
i\leq 8$
\end{itemize}
définis par les formules:
\begin{tabbing}
\quad\= (1) \= $(m_1, m_2, n)$ \= $\longmapsto$ \= $(m_1+m_2, n)$\\
\> (2) \> $(m_1, m_2, n)$ \> $\longmapsto$ \> $(m_1, n)$\\
\> (3) \> $(m_1, m_2, n)$ \> $\longmapsto$ \> $(m_2, n)$\\
\> (4) \> $(m, n_1, n_2)$ \> $\longmapsto$ \> $(m, n_1+n_2)$\\
\> (5) \> $(m, n_1, n_2)$ \> $\longmapsto$ \> $(m, n_1)$\\
\> (6) \> $(m, n_1, n_2)$ \> $\longmapsto$ \> $(m, n_2)$\\
\> (7) \> $(m, a, n)$ \> $\longmapsto$ \> $(ma, n)$\\
\> (8) \> $(m, a, n)$ \> $\longmapsto$ \> $(m, an)$;
\end{tabbing}
où la formule (i) décrit l'application (i)(X) pour les objets
variables $X$ de $E$. Notons $L$ le foncteur « faisceau abélien
libre engendré ». Il est clair que le plus grand quotient (au sens
des faisceaux abéliens) de $L(M\times N)$ qui égalise le morphisme
$L(1)$ à $L(2)+L(3)$, $L(4)$ à $L(5)+L(6)$ et $L(7)$ à $L(8)$
représente le foncteur des morphismes $A$-bilinéaires de $M\times N$
dans $P$.

\setcounter{subsection}{7}





\subsection{}\etiquette[12.8]{IV.subsec12.8}
On constate que le foncteur $P\mapsto \Hom_A(N,\scrHom_{\sisi{\underline{\underline{\ZZ}}}{\underline{\ZZ}}}(M,P))$ est
aussi canoniquement isomorphe au foncteur des applications
$A$-bilinéaires de $M\times N$ dans $P$. Par suite le produit
tensoriel $M\otimes_A N$ représente aussi le foncteur
$\Hom_A(N,\scrHom_{\sisi{\underline{\underline{\ZZ}}}{\underline{\ZZ}}}(M,P))$. On a donc des isomorphismes, fonctoriel
en tout les arguments:
\begin{align*}%\etiquette[12.8.1]{IV.eq12.8.1}
\Hom_Z(M\otimes_A N,P) &\simeq \Hom_A(M,\scrHom_{\sisi{\underline{\underline{\ZZ}}}{\underline{\ZZ}}}(N,P)){\quoi},\tag{12.8.1}\\
\Hom_Z(M\otimes_A N,P) &\simeq \Hom_A (N,\scrHom_{\sisi{\underline{\underline{\ZZ}}}{\underline{\ZZ}}}(M,P)){\quoi}.\tag{12.8.2}
\end{align*}

\subsection{}\etiquette[12.9]{IV.subsec12.9}
Il\pageoriginale résulte de (12.8), ou bien de
(12.8.1) et de (\hyperref[IV.subsec12.2]{12.2}) que le foncteur
$\otimes_A$ commute aux limites inductives en ses deux arguments.

\begin{enonce*}{Proposition 12.10}\etiquette[12.10]{IV.prop12.10}
Soient $C$ un $𝒰$-site, $A'$ un préfaisceau d'anneaux sur
$C$, $M'$ (\resp $N'$) un $A'$-module à droite (\resp à gauche),
$A$, $M$, $N$ les faisceaux associés à $A'$, $M'$ et $N'$
respectivement. Le faisceau associé au préfaisceau $X\mapsto
M'(X)\otimes_{A(X)}N'(X)$ $(X\in \ob C)$ est canoniquement isomorphe
au faisceau $M\otimes_A N$.
\end{enonce*}

Le préfaisceau $X\mapsto M'(X)\otimes_{A'(X)}N'(X)$ peut se
construire à partir des préfaisceaux $M'$, $N'$ et $A'$ par les
opérations indiquées dans la démonstration de \hyperref[IV.prop12.7]{12.7}.
Comme le foncteur « faisceau associé » commute aux limites
projectives et inductives finies et aux foncteurs « objet abélien
libre engendré », la formation du produit tensoriel commute au
foncteur « faisceau associé ».

\begin{enonce*}{Proposition 12.11}\etiquette[12.11]{IV.prop12.11}
Soient $(E,A)$ un topos annelé, $M$ un $A$-Module à droite, $N$ un
$A$-Module à gauche, $X$ un objet de $E$.
\begin{enumerate}
\renewcommand{\theenumi}{\alph{enumi}}
\renewcommand{\labelenumi}{\theenumi)}
\item On a un isomorphisme canonique
$$
j^{\ast}_X(M\otimes_A N)\simeq (j^{\ast}_XM)\otimes_A
(j^{\ast}_XN){\quoi}.
$$
Soient de plus $P$ un $A/X$-Module à droite et $Q$ un $A/X$ module à
gauche.

\item On a des isomorphismes canoniques (formules de projection):
\begin{align*}
j_{X!}(P\otimes_{A| X} j^{\ast}_XN) &\simeq (j_{X!}P)\otimes_A N\\
j_{X!}(j^{\ast}_XM\otimes_{A| X}Q) &\simeq M\otimes_A
(j_{X!}Q){\quoi}.
\end{align*}
\end{enumerate}
\end{enonce*}

\begin{itemize}
\item[Démontrons a).] Pour tout faisceau abélien $R$ sur $E_{/X}$, on
a la suite d'isomorphismes fonctoriels en $R$:
\begin{alignat*}{2}
&\Hom_{\sisi{\underline{\underline{\ZZ}}}{\underline{\ZZ}}}(j^{\ast}_X(M\otimes_A N),R)\simeq \Hom_{\sisi{\underline{\underline{\ZZ}}}{\underline{\ZZ}}} (M\otimes_A N,
j_{X\ast}R) &\quad& \eqhyperref[IV.eq11.2.2.1]{11.2.2.1}\\
&\Hom_{\sisi{\underline{\underline{\ZZ}}}{\underline{\ZZ}}}(M\otimes_A N, j_{X\ast}R)\simeq \Hom_A(M, \scrHom_Z(N,
j_{X\ast}R)) && (12.8.1)\\
&\Hom_A (M, \scrHom_Z(N, j_{X\ast}R))\simeq
\Hom_A(M,j_{X\ast}\scrHom_{\sisi{\underline{\underline{\ZZ}}}{\underline{\ZZ}}}(j^{\ast}_XN,R)) && (\hyperref[IV.prop12.3]{12.3})\\
& \Hom_A(M, j_{X\ast}\scrHom_{\sisi{\underline{\underline{\ZZ}}}{\underline{\ZZ}}}(j^{\ast}_XN,R))\simeq \Hom_{A|
X}(j^{\ast}_XM, \scrHom_{\sisi{\underline{\underline{\ZZ}}}{\underline{\ZZ}}} (j^{\ast}_XN,R)) && \eqhyperref[IV.eq11.2.2.1]{11.2.2.1}\\
&\Hom_{A/X}(j^{\ast}_XM, \scrHom_{\sisi{\underline{\underline{\ZZ}}}{\underline{\ZZ}}} (j^{\ast}_XN,R)\simeq \Hom_{\sisi{\underline{\underline{\ZZ}}}{\underline{\ZZ}}}
(j^{\ast}_XM\otimes_A j^{\ast}_X N,R) && (12.8.1)
\end{alignat*}

Exhibons\pageoriginale le premier isomorphisme de b). Pour tout
faisceau abélien $R$, on a une suite d'isomorphismes fonctoriels en
$R$:
\begin{alignat*}{2}
& \Hom_Z(j_{X!}(P\otimes_{A| X}j^{\ast}_XN),R)\simeq
\Hom_Z(P\otimes_{A| X} j^{\ast}_X N, j^{\ast}_X R) &\quad&
(\hyperref[IV.prop11.3.1]{11.3.1})\\
& \Hom_Z(P\otimes_{A| X}j^{\ast}_XN, j^{\ast}_XR)\simeq
\Hom_{A/X}(P,\scrHom_Z(j^{\ast}_XN, j^{\ast}_XR)) && (12.8.1)\\
&\Hom_{A| X}(P,\scrHom_Z (j^{\ast}_XN, j^{\ast}_XR))\simeq
\Hom_{A| X}(P, j^{\ast}_X\scrHom_Z (N,R)) && (\hyperref[IV.prop12.3]{12.3})\\
&\Hom_{A| X}(P, j^{\ast}_X\scrHom_Z(N,R))\simeq \Hom_A (j_{X!}P,
\scrHom_Z (N,R)) && (\hyperref[IV.prop11.3.1]{11.3.1})\\
&\Hom_A (j_{X!}P, \scrHom_Z(N,R))\simeq \Hom_Z (j_{X!}P\otimes_A
N,R) && (12.8.1)
\end{alignat*}

Le deuxième isomorphisme s'obtient de manière analogue à l'aide de
(12.8.2).
\end{itemize}

\begin{enonce*}{Corollaire 12.12}\etiquette[12.12]{IV.cor12.12}
Soient $E$ un topos, $A$, $B$, $C$ trois faisceaux d'anneaux sur
$E$, $M$ un $B-A$ biModule, $N$ un $A-C$ biModule. Le faisceau
abélien $M\otimes_A N$ est muni canoniquement d'une structure de
$B-C$ biModule. En particulier, lorsque $A$ est un faisceau
d'anneaux commutatifs et lorsque $M$ et $N$ sont des $A$-Modules,
$M\otimes_AN$ est muni canoniquement d'une structure de $A$-Module.
Les isomorphismes de $\hyperref[IV.subsec12.9]{12.9}$ sont des isomorphismes de
biModules.
\end{enonce*}

Pour tout objet $X$ de $E$, $j^{\ast}_XM$ est un $A| X$-Module à
droite sur lequel opère, à gauche, l'anneau $B(X)$ et de même,
$j^{\ast}_XN$ est un $A/X$-Module à gauche sur lequel opère, à
droite, $C(X)$. Par fonctorialité, le faisceau abélien
$j^{\ast}_XM\otimes_{A/X} j^{\ast}_XN\simeq j^{\ast}_X (M\otimes_A
N)$ est muni d'une structure de $B(X)-C(X)$ « objet », structure
qui varie fonctoriellement en $X$. Par suite $M\otimes_A N$ est muni
d'une structure de $B-C$ biModule. Cette structure de $B(X)-C(X)$
objet sur $j^{\ast}_X(M\otimes_A N)$ se reflète de façon évidente
sur le foncteur « morphisme $A$-bilinéaire ». On constate alors
que, lorsque $A$ est commutatif et lorsque $M$ et $N$ sont des
$A$-Modules, les structures de $A$-Modules à droite et à gauche
qu'on obtient sont « égales »  et par suite $M\otimes_A N$ est
dans ce cas un $A$-Module. La dernière assertion est laissée au
lecteur.

\begin{enonce*}{Corollaire 12.13}\etiquette[12.13]{IV.cor12.13}
Soient $(E,A)$ un topos annelé, $M$ un $A$-Module à
gauche\pageoriginale (\resp à droite), $X$ un objet de $E$. On a
(avec la notation $A_X$ de $\hyperref[IV.prop11.3.3]{11.3.3}$) un isomorphisme
canonique
$$
A_X\otimes_AM\simeq J_{X!} j^{\ast}_XM\text{ (\resp } M\otimes_A
A_X\simeq j_{X!}j^{\ast}_XM)){\quoi}.
$$
\end{enonce*}

\noindent résulte des formules de projections (\hyperref[IV.prop12.11]{12.11}).

\begin{enonce*}{Proposition 12.14}\etiquette[12.14]{IV.prop12.14}
Soient $E$ un topos, $A$ et $B$ deux faisceaux d'anneaux sur $E$,
$M$ un $A$-Module à droite, $N$ un $A-B$ biModule, $P$ un $B$-Module
à droite. On a un isomorphisme canonique
$$
\Hom_B(M\otimes_A N,P)\simeq \Hom_A (M,\scrHom_B (N,P)){\quoi}.
$$
\end{enonce*}

\noindent De (12.8.1) on tire un morphisme $A$-bilinéaire
canonique $\scrHom_Z(N,P)\times N\rightarrow P$; d'où, en se
restreignant à $\scrHom_B(N,P)\times N$ (qui est un sous-faisceau),
un morphisme $A$-bilinéaire $\scrHom_B(N,P)\times N\rightarrow P$
dont on vérifie immédiatement qu'il est $B$-linéaire sur
le deuxième
facteur. On a donc un morphisme canonique de $B$-Modules
$\scrHom_B(N,P)\otimes_A N\rightarrow P$; d'où une application
canonique, fonctorielle en $M$:
\begin{equation*}
\Hom_A(M,\scrHom_B(N,P))\longrightarrow \Hom_B (M\otimes_A
N,P){\quoi}.\tag{12.14.1}
\end{equation*}\etiquette[12.14.1]{IV.eq12.14.1}
Montrons que ce morphisme de foncteurs est un isomorphisme. Comme
les deux membres transforment les limites inductives de $M$ en
limites projectives, il suffit de montrer que \eqhyperref[IV.eq12.14.1]{12.14.1} est un
isomorphisme, lorsque $M$ parcourt une famille génératrice de la
catégorie $E_A$. Il suffit donc de montrer que \eqhyperref[IV.eq12.14.1]{12.14.1} est
un isomorphisme lorsque $M=A_X$ où $X$ est un objet de $E$. La
vérification est alors immédiate.


\section{Morphisme de topos annelés}\etiquette[13]{IV.sec13}

\begin{enonce*}{Définition 13.1}\etiquette[13.1]{IV.def13.1}
Soient $(E,A)$ et $(E',A')$ deux topos annelés. Un morphisme de
topos annelé $u:(E,A)\rightarrow (E',A')$ est un couple $(m,\theta)$
où $m:E\rightarrow E'$ est un morphisme de topos {\rm
(\hyperref[IV.def3.1]{3.1})} et $\theta:m^{\ast}A'\rightarrow A$ est un
morphisme d'Anneaux.
\end{enonce*}

\setcounter{subsection}{1}
\subsubsection{}\etiquette[13.1.1]{IV.subsubsec13.1.1}
Comme le foncteur $m^{\ast}:E'\rightarrow E$ est adjoint à gauche au
foncteur $m_{\ast}:E\rightarrow E'$, se donner un morphisme de topos
annelés $(m,\theta):(E,A)\rightarrow (E',A')$ revient à se donner un
morphisme de topos $m:E\rightarrow E'$ et un morphisme de faisceaux
d'anneaux $\Theta:A'\rightarrow m_{\ast}A$.

\subsection{}\etiquette[13.2]{IV.subsec13.2}
A\pageoriginale un morphisme de topos annelés
$u=(m,\theta):(E,A)\rightarrow (E',A^{\ast})$, on associe deux
foncteurs remarquables entre les catégories de Modules:

\subsubsection{Le foncteur image directe pour les
  Modules:}\etiquette[13.2.1]{IV.subsubsec13.2.1}
Soit $M$ un $A$-Module à gauche (\resp à droite). L'objet
$m_{\ast}M$ est muni canoniquement d'une structure de
$m_{\ast}A$-Module; d'où par restriction des scalaires par le
morphisme canonique $\Theta':A'\rightarrow m_{\ast}A$, un
$A'$-Module noté $u_{\ast}(M)$ et appelé \emph{l'image directe de
$M$ par le morphisme $u$.}

\subsubsection{Le foncteur image réciproque pour les
  Modules:}\etiquette[13.2.2]{IV.subsubsec13.2.2}
Soit $N$ un $A'$-Module à gauche (\resp à droite). L'objet
$m^{\ast}N$ de $E$ est muni canoniquement d'une structure de
$m^{\ast}A'$-Module à gauche (\resp à droite). Le $A$-Module à
gauche $A\otimes_m\ast_{A'}m^{\ast}N$ (\resp à droite
$m^{\ast}N\otimes_m\ast_{A'}A$) où $A$ est muni de la structure de
$m^{\ast}A'$-Module définie par $\theta:m^{\ast}A'\rightarrow A$,
est noté $u^{\ast}N$ et est appelé l'\emph{image réciproque du
Module $N$ par le morphisme de topos annelé} $u$.

\subsubsection{}\etiquette[13.2.3]{IV.subsubsec13.2.3}
On notera que pour tout $A$-Module $M$, le faisceau d'ensembles et
le faisceau abélien sous-jacent à $u_{\ast}M$ est le faisceau
$m_{\ast}M$. Aussi emploie-t-on le plus souvent la notation
$u_{\ast}$ pour désigner le foncteur $m_{\ast}$: image directe pour
les faisceaux d'ensembles. En revanche pour un $A'$-Module $N$, le
faisceau d'ensembles ou le faisceau abélien sous-jacent à
$u^{\ast}N$ n'est ni égal ni isomorphe en général à $m^{\ast}N$
(sauf toutefois lorsque $\theta$ est un isomorphisme). Il y a donc
lieu de distinguer entre l'image réciproque pour les Modules et
l'image réciproque pour les faisceaux d'ensembles ou les faisceaux
abéliens. On utilise le plus souvent la notation
$u^{-1}:E'\rightarrow E$ pour désigner le foncteur $m^{\ast}$ image
inverse pour les faisceaux d'ensembles. Le foncteur $u^{-1}$ est
appelé le foncteur \emph{image réciproque « ensembliste »} par
le morphisme de topos annelé $u$, par opposition avec l'image
réciproque « au sens modules »  $u^{\ast}$.

\begin{enonce*}{Définition 13.3}\etiquette[13.3]{IV.def13.3}
Soient $(C,A)$ et $(C',A')$ deux $𝒰$-sites annelés. Un
morphisme de sites annelés $u:(C,A)\rightarrow (C',A')$ est un
couple $(m,\theta)$ où $m$ est un morphisme du site $C$ dans le site
$C'$ (\hyperref[IV.subsec4.9]{4.9}) et $\theta:m^{\ast}A'\rightarrow
A$\pageoriginale un morphisme d'Anneaux.
\end{enonce*}

\setcounter{subsection}{3} \setcounter{subsubsection}{0}
\subsubsection{}\etiquette[13.3.1]{IV.subsubsec13.3.1}
Un morphisme de topos annelé est un morphisme de $𝒰$-sites
annelés. Un morphisme de sites annelés donne naissance à un
morphisme entre les topos annelés correspondants
(\hyperref[IV.subsubsec11.1.1]{11.1.1} et \hyperref[IV.subsubsec4.9.1]{4.9.1}).

\begin{enonce*}{Proposition 13.4}\etiquette[13.4]{IV.prop13.4}
Soit $u:(E,A)\rightarrow (E',A')$ un morphisme de topos annelés.
\begin{enumerate}
\renewcommand{\theenumi}{\alph{enumi}}
\renewcommand{\labelenumi}{\theenumi)}
\item Pour un $A$-Module à gauche (\resp à droite) variable $M$, et
pour un $A'$-Module à gauche (\resp à droite) variable $N$, on a des
isomorphismes bifonctoriels canoniques (dits isomorphismes
d'adjonction):
\begin{align*}
\Hom_A(u^{\ast}N,M) &\simeq
\Hom_{A'}(N,u_{\ast}M){\quoi},\tag{13.4.1}\\
u_{\ast}\scrHom_A(u^{\ast}N,M) &\simeq
\scrHom_{A'}(N,u_{\ast}M){\quoi}.\tag{13.4.2}
\end{align*}

\item Pour un objet variable $X$ de $E'$, on a un isomorphisme
canonique, fonctoriel en $X$:
\begin{equation*}
u^{\ast}A'_X\longrightarrow A_{u^{-1}(X)}{\quoi},\tag{13.4.3}
\end{equation*}\etiquette[13.4.3]{IV.eq13.4.3}
où $A'_X$ et $A_{u^{-1}(X)}$ désignent les Modules libres engendrés
$(\hyperref[IV.prop11.3.3]{11.3.3})$.

\item Lorsque $A'$ est commutatif et lorsque le morphisme canonique
$u^{-1}A'\rightarrow A$ est central (\resp lorsque le morphisme
canonique $u^{-1}A'\rightarrow A$ est un isomorphisme) on a, pour un
$A'$-Module à droite $M$ variable et un $A'$-Module à gauche $N$
variable, un isomorphisme bifonctoriel canonique:
\begin{align*}
u^{\ast}M\otimes_A u^{\ast}N &\simeq
u^{\ast}(M\otimes_{A'}N)\tag{13.4.4}\\
u^{\ast}M\otimes_A u^{\ast}N &\simeq
u^{-1}(M\otimes_{A'}N)){\quoi}.\tag*{(\resp (13.4.5)}
\end{align*}
\end{enumerate}
\end{enonce*}

On a tout d'abord un isomorphisme canonique
$\Hom_A(u^{\ast}N,M)\simeq \Hom_{u^{-1}A'}(u^{-1}N,M)$
(\sisi{\hyperref[IV.cor12.12]{12.12}}{\hyperref[IV.prop12.14]{12.14}}), puis un isomorphisme
$\Hom_{u^{-1}A'}(u^{-1}N,M)\simeq \Hom_{A'}(N,u_{\ast}M)$
(\hyperrefext[III][III.prop1.7]{1.7}); d'où (13.4.1). Exhibons l'isomorphisme
(13.4.2). Pour tout objet $X$ de $E'$, on a la suite
d'isomorphismes fonctoriels:
\begin{alignat*}{3}
& \Hom_{E'}(X,u_{\ast}\scrHom_A (u^{\ast}N,M))\simeq && \Hom(
u^{-1}X,
 \scrHom_A (u^{\ast}N, M)) &\quad& \text{(adjonction){\quoi},}\\
&\Hom_A(u^{\ast}N,\scrHom_E (u^{-1}X, M)) \simeq &&\qquad ''\qquad ''\qquad
''\qquad '' && (\hyperref[IV.prop12.1]{12.1}){\quoi},\\
& \Hom_A (u^{\ast}N,\scrHom_E(u^{-1}X, M))\simeq &&\Hom_{A'}(N,
u_{\ast} \scrHom_E (u^{-1} X, M)) && (13.4.1){\quoi},\\
& \Hom_{A'}(N,\scrHom_{E'}(X,u_{\ast}M))\simeq &&\qquad ''\qquad  ''\qquad
 ''\qquad '' && \eqhyperref[IV.eq10.3.1]{10.3.1}{\quoi},\\
& \qquad ''\qquad ''\qquad ''\qquad '' \simeq &&
 \Hom_{A'}(X,\scrHom_{A'}(N,u_{\ast}M)) && (\hyperref[IV.prop12.1]{12.1}){\quoi}.
\end{alignat*}
La\pageoriginale formule (13.4.5) s'obtient alors à partir de
la formule 13.3.2. par adjonction (définition du produit tensoriel
(\hyperref[IV.prop12.7]{12.7})). La formule (13.4.4) se déduit de
(13.4.5) en utilisant la commutativité du produit tensoriel
lorsque l'anneau de base est commutatif (\hyperref[IV.subsec12.8]{12.8}).
Enfin, pour démontrer \eqhyperref[IV.eq13.4.3]{13.4.3}, on considère la suite
d'isomorphismes
\begin{alignat*}{3}
& \Hom_A(u^{\ast}A'_X, M) && \simeq\Hom_{A'}(A'_X, u_{\ast}M)
&\quad&
(13.4.1){\quoi},\\
& \Hom_{E'}(X, u_{\ast}M) && \simeq\quad ''\qquad ''\qquad ''&&
(\hyperref[IV.prop11.3.3]{11.3.3}){\quoi},\\
&\quad ''\qquad ''\qquad ''&& \simeq\Hom_E(u^{-1}X, M) &&
\text{(adjonction)}{\quoi},\\
&\Hom_A(A_{u^{-1}X},M) &&\simeq\quad ''\qquad ''\qquad ''&&
(\hyperref[IV.prop11.3.3]{11.3.3}){\quoi}.
\end{alignat*}

\begin{enonce*}{Corollaire 13.5}\etiquette[13.5]{IV.cor13.5}
Soient $(E,A)$ un topos annelé, $x:P\rightarrow E$ un point de $E$,
$F\mapsto F_x$ le foncteur fibre associé (\hyperref[IV.def6.1]{6.1}). Le
foncteur fibre en $x$ transforme le produit tensoriel des
$A$-Modules en produit tensoriel (ordinaire) des $A_x$-modules.
\end{enonce*}

Le produit tensoriel dans $P$ est le produit tensoriel ordinaire des
modules ($P$ le topos ponctuel est la catégorie des ensembles).
L'assertion résulte donc de (13.4.4).

\begin{enonce*}{Corollaire 13.6}\etiquette[13.6]{IV.cor13.6}
Soit $u:(E,A)\rightarrow (E',A')$ un morphisme de topos annelés. Le
foncteur $u_{\ast}$ image directe pour les Modules à droite ou à
gauche commute aux limites projectives et en particulier est exact à
gauche. Le foncteur $u^{\ast}$ image réciproque pour les Modules à
droite ou à gauche commute aux limites inductives et en particulier
est exact à droite.
\end{enonce*}

Ceci résulte de la formule (13.4.1) (\hyperrefext[I][I.prop2.11]{2.11}).

\setcounter{subsection}{6}
\subsection{}\etiquette[13.7]{IV.subsec13.7}
On notera que le foncteur $u^{\ast}$, image réciproque pour les
Modules, n'est pas, en général, exact, alors que le foncteur
$u^{-1}$, image réciproque ensembliste, est exact. On peut cependant
affirmer l'exactitude de $u^{\ast}$ lorsque le morphisme canonique
$u^{-1}A'\rightarrow A$ est plat (à droite\pageoriginale ou à
gauche) (V 1.8). C'est en particulier le cas lorsque le morphisme
canonique $u^{-1}A'\rightarrow A$ est un isomorphisme.

\subsection{}\etiquette[13.8]{IV.subsec13.8}
Soient $C$ un $𝒰$-site, $A'$ un préfaisceau d'anneaux sur
$C$, et notons $A$ le faisceau associé à $A'$. La catégorie des
préfaisceaux de $A'$-modules qui sont des faisceaux est équivalente
à la catégorie des $A$-Modules. On utilise parfois la notation
$\Hom_{A'}(M,N)$ pour désigner le groupe $\Hom_A(M,N)$
(\hyperref[IV.subsubsec11.1.2]{11.1.2}). De même, et abusivement, on utilise les
notations $\scrHom_{A'}(M,N)$ et $M\otimes_{A'}N$ pour désigner les
faisceaux $\scrHom_A(M,N)$ et $M\otimes_A N$. Lorsque $A'=k_E$ est
le préfaisceau constant, défini par un anneau ordinaire $k$, on
écrit aussi $\scrHom_k$, $\otimes_k$ au lieu de $\scrHom_{k_E}$,
$\otimes_{k_E}$

\begin{enonce*}[remark]{Exercice 13.9}\etiquette[13.9]{IV.exer13.9}
\emph{Topos localement annelés} (\cf [\hyperref[IV.sec9]{9}] pour plus de
renseignements dans l'ordre d'idées qui suit).

Soit $(E,{\Oo}_E)$ un topos commutativement annelé. Pour $X\in
\ob E$ et $f\in {\Oo}_E(X)$, soit $X_f$ le plus grand sous-objet
de $X$ sur lequel $f$ soit inversible.
\begin{enumerate}
\renewcommand{\theenumi}{\alph{enumi}}
\renewcommand{\labelenumi}{\theenumi)}
\item Montrer que pour $X'\in \ob E_{/X}$, $f_{X'}$ est inversible si
et seulement si le morphisme structural $X'\rightarrow X$ se
factorise par $X_f$, et que pour $f$, $g\in {\Oo}_E(X)$, on a
$$
X_{fg}=X_f\cap X_g{\quoi}.
$$

\item Montrer que les conditions suivantes (i) à (iii) sont
équivalentes:
\begin{enumerate}
\renewcommand{\theenumii}{\roman{enumii}}
\renewcommand{\labelenumii}{(\theenumii)}
\item Pour $X\in \ob E$ et $f$, $g\in {\Oo}_E(X)$, on a
$$
X_{f+g}\subset \Sup(X_f,X_g)
$$

\item Pour $X\in \ob E$ et $f$, $g\in {\Oo}_E(X)$ tels que $f+g$ soit
inversible, on a $X=\Sup(X_f,X_g)$.

\item Pour $X\in \ob E$ et $f\in {\Oo}_E(X)$, on a
$$
X=\Sup(X_f, X_{1-f}){\quoi}.
$$

Montrer que ces conditions impliquent la condition suivante (iv),
et\pageoriginale sont équivalentes à cette dernière si $E$ a
suffisamment de points:

\item Pour tout point $p$ de $E$, l'anneau fibre ${\Oo}_{E,p}$ est un
\emph{anneau local}.
\end{enumerate}

Lorsque les conditions équivalentes (i) à (iii) sont satisfaites, on
dira que $(E,{\Oo}_E)$ est un \emph{topos localement annelé}, et
on dit de même qu'un site annelé est un \emph{site localement
annelé} si le topos annelé correspondant est un topos localement
annelé.

\item Soient $(E,{\Oo}_E)$ et $(E',{\Oo}_{E'})$ deux topos localement
annelés, et $f$ un morphisme de topos annelés du premier dans le
second. Montrer que la condition (i) suivante implique la condition
(ii) et lui est équivalente si $E$ a suffisamment de points:
\begin{enumerate}
\renewcommand{\theenumii}{\roman{enumii}}
\renewcommand{\labelenumii}{(\theenumii)}
\item Pour tout $X\in \ob E'$ et $s'\in {\Oo}_{E'}(X')$, posant
$X=f^{-1}(X)$, $s=f^{\ast}(s')$, on~a
$$
X'_{s'}=f^{-1}(X_s){\quoi}.
$$

\item Pour tout point $p$ de $E$, posant $p'=f(p)$, l'homomorphisme
naturel sur les fibres
$$
{\Oo}_{E',p'}\longrightarrow {\Oo}_{E,p}
$$
est un homomorphisme \emph{local} d'anneaux locaux.
\end{enumerate}

Lorsque la condition (i) est satisfaite, on dira que $f$ est un
\emph{morphisme de topos localement annelés}, ou encore un
\emph{morphisme admissible de topos localement annelés} si une
confusion est à craindre. On désigne par $\Hhomtoplocan(E,E')$ la
sous-catégorie pleine de la catégorie $\Hhomtopan(E,E')$ de tous les
morphismes de topos annelés de $E$ dans $E'$ définie par les
morphismes admissibles de $E$ dans $E'$.

\item Supposons que $(E,{\Oo}_E)$ soit le topos annelé
  associé à un
espace topologique annelé $(X,{\Oo}_X)$ (\hyperref[IV.subsec2.1]{2.1}).
Montrer que pour que $(E,{\Oo}_E)$ soit localement annelé, il
faut et il suffit que $(X,{\Oo}_X)$ soit localement annelé, \ie
que pour tout $x\in X$, la fibre ${\Oo}_{X,x}$ soit un anneau
local. (Noter que dans le critère (iv) de b), il suffit de prendre
le point $p$ dans une famille \emph{conservatrice} de points de
$E$). Soit\pageoriginale $f:(X,{\Oo}_X)\rightarrow (X',
{\Oo}_{X'})$ un morphisme d'espaces annelés, où ${\Oo}_X$ et
${\Oo}_{X'}$ sont des faisceaux d'anneaux locaux. Montrer que le
morphisme de topos annelés correspondant
$\Top(X,{\Oo}_X)\rightarrow \Top(X', {\Oo}_{X'})$ est
admissible si et seulement si il en est de même pour le morphisme
$f$, \ie; si et seulement si pour tout $x\in X$,
${\Oo}_{X',f(x)}\rightarrow {\Oo}_{X,x}$ est un
homomorphisme \emph{local} d'anneaux locaux.

\item Soient $(E,{\Oo}_E)$ et $(E',{\Oo}_{E'})$ deux topos localement
annelés, tels que $E$ ait suffisamment de points et que
$(E',{\Oo}_{E'})$ soit équivalent au topos annelé défini par un
\emph{schéma} $(X', {\Oo}_{X'})$; prouver que la catégorie
$\Hhomtoplocan(E,E')$ est équivalente à une catégorie
\emph{discrète}, \ie que c'est un groupoïde (tout morphisme est un
isomorphisme) rigide (les groupes d'automorphismes des objets sont
les groupes unité).

\noindent (Hint: se ramener au cas où $(E,{\Oo}_E)$ est le topos
ponctuel annelé par un corps). On se rappellera que, par contre,
$\Hhomtop(E,E')$ n'est pas en général équivalent à une
catégorie
discrète, même si $E$ et $E'$ sont des topos définis par des
schémas
(et même si $E$ est le topos ponctuel), \cf
\hyperref[IV.subsubsec4.2.3]{4.2.3}.

\item Soient $E$ un topos localement annelé, $P$ le topos annelé
défini par l'espace annelé $\Spec k$, où $k$ est un corps. Montrer
que les morphismes admissibles de topos localement annelés
$P\rightarrow E$ correspondent aux couples $(p,u)$ d'un point $p$ de
$E$, et d'une injection de $k(p)$ dans $k$. où $k(p)$ est le corps
résiduel de l'anneau local ${\Oo}_{E,p}$. Généraliser en un
énoncé exhibant les morphismes admissibles d'un topos localement
annelé ponctuel dans $E$ (généralisant EGA 1 2.4.4).

On appelle \emph{point géométrique d'un topos localement
  annelé} $E$
tout morphisme admissible dans $E$ du topos localement annelé défini
par un espace annelé de la forme $\Spec(k)$, où $k$ est un corps
\emph{algébriquement clos}. Définir la \emph{catégorie des points
géométriques de} $E$ (sous-entendu: correspondants à des corps $k\in
𝒰$), notée $\Pptgeom(E)$, et un foncteur canonique
$\Pptgeom(E)\rightarrow \Ppoint(E)$. Montrer que ce foncteur est
fidèle si\pageoriginale et seulement si $E$ n'a pas de point \ie
$\Ppoint(E)$ est la catégorie vide.

\item Soient $E$ un $𝒰$-topos, $V$ un univers tel que $U\in
V$. Définir une 2-catégorie
$(𝒱\text{-}𝒰\text{-}\Top)_{/E}$ en termes de l'objet
$E$ de la 2-catégorie $(𝒱\text{-}𝒰\text{-}\Top)$
(\hyperref[IV.subsubsec3.3.1]{3.3.1}), en s'inspirant de \hyperref[IV.exer5.14]{5.14} a).
Lorsque $E$ est annelé par un Anneau $A$, définir de même une
2-catégorie $(V\text{-}𝒰\text{-}\text{Topan})_{/E}$, et
lorsque $E$ est localement annelé, utilisant la notion de morphisme
admissible (\cf c)), définir une 2-catégorie
$(𝒱\text{-}𝒰\text{-}\text{Topanloc})_{/E}$,
admettant $\Pptgeom(E)$ (\cf f)) comme sous-catégorie pleine.
\end{enumerate}
\end{enonce*}

\section{Modules sur un topos défini par recollement}\etiquette[14]{IV.sec14}

\subsection{}\etiquette[14.1]{IV.subsec14.1}
Soient $(E,A)$ un topos annelé, $U$ un sous-topos ouvert de $E$, $F$
le sous-topos fermé complémentaire, $j:U\rightarrow E$ et
$i:F\rightarrow E$ les morphismes canoniques (\hyperref[IV.subsec9.3]{9.3}).
Le topos $U$ est annelé par $j^{\ast}A$ (\sisi{\hyperref[IV.subsubsec11.2.2]{11.2.2}}{\hyperref[IV.subsubsec11.2.1]{11.2.1}})
et le morphisme de topos $i$, complété de manière évidente, devient
un morphisme de topos annelés. Dans ce numéro, le topos $F$ sera
annelé par le faisceau $i^{\ast}A$, noté $A| F$, et le morphisme
$j$, complété de manière évidente, est alors un morphisme de topos
annelés.


\subsection{}\etiquette[14.2]{IV.subsec14.2}
On a donc deux morphismes de topos annelés:
$$
j:(U,A| U)\longrightarrow (E,A)\text{ et } i:(F,A| F)\longrightarrow
(E,A){\quoi};
$$
d'où cinq foncteurs entre les catégories de Modules (à gauche pour
fixer les idées) correspondantes:
\[\vcenter{
\xymatrix@C=2cm@R=1.5cm@M=.25cm{ ({}_{A|U}U) \ar@<14pt>[r]^{j_!}
\ar@<-14pt>[r]^{j_{\ast}}&
  ({}_{A}E) \ar[l]_{j^{\ast}} \ar@<7pt>[r]^{i^{\ast}} &
  ({}_{A|F}F)\quad. \ar@<7pt>[l]_{i_{\ast}}
}}\leqno(14.2.1)
\]\etiquette[14.2.1]{IV.eq14.2.1}
Chaque foncteur du diagramme \eqhyperref[IV.eq14.2.1]{14.2.1} est adjoint à gauche à
celui qui se trouve au-dessous de lui.

\subsection{}\etiquette[14.3]{IV.subsec14.3}
On\pageoriginale sait que le foncteur $X\mapsto (j^{\ast}X,
i^{\ast}X; i^{\ast}X\rightarrow i^{\ast}j_{\ast}j^{\ast}X)$ est une
équivalence de $E$ dans le topos $(U,F, i^{\ast}j_{\ast})$
(\hyperref[IV.theo9.5.4]{9.5.4}). Cette équivalence induit une équivalence entre
les catégories de Modules correspondantes et par suite le foncteur
$P\mapsto (j^{\ast}P, i^{\ast}P; i^{\ast}P\rightarrow
i^{\ast}j_{\ast}j^{\ast}P)$ est une équivalence, notée $\Phi$, de la
catégorie ${}_AE$ dans la catégorie $({}_{A| U}U, {}_{A| F}F,
i^{\ast}j_{\ast})$. Les foncteurs de (14.2.1) composés avec $\Phi$
ou $\Phi^{-1}$ sont alors les foncteurs:
\begin{align*}
\Phi\circ j_{\ast} &: M\mapsto (M,i^{\ast} j_{\ast}M;
i^{\ast}j_{\ast}M\xrightarrow{\id} i^{\ast}j_{\ast}M){\quoi},\\
j^{\ast}\circ \Phi^{-1} &: (M,N; N\rightarrow
i^{\ast}j_{\ast}M)\mapsto M{\quoi},\\
\Phi\circ j_{!} &: M\mapsto (M,0;0\rightarrow i^{\ast}
j_{\ast}M){\quoi},\\
\Phi\circ i_{\ast} &: N\mapsto (0,N;N\rightarrow 0){\quoi},\\
i^{\ast}\circ \Phi^{-1} &: (M,N;N\rightarrow
i^{\ast}j_{\ast}M)\mapsto N{\quoi}.
\end{align*}
Le lecteur pourra, à titre d'exercice, expliciter les différents
morphismes d'adjonction.

\subsection{}\etiquette[14.4]{IV.subsec14.4}
Notons
\begin{equation*}
i^{!}:{}_AE\longrightarrow {}_{A| F}F\tag{14.4.1}
\end{equation*}\etiquette[14.4.1]{IV.eq14.4.1}
le foncteur défini par la formule:
\begin{equation*}
i^{!}\circ
\Phi^{-1}(M,N;N\xrightarrow{u}i^{\ast}j_{\ast}M)=\Ker(u){\quoi}.\tag{14.4.2}
\end{equation*}\etiquette[14.4.2]{IV.eq14.4.2}

\begin{enonce*}{Proposition 14.5}\etiquette[14.5]{IV.prop14.5}
Le foncteur $i^{!}$ est adjoint à droite au foncteur $i_{\ast}$. Le
morphisme d'adjonction $i_{\ast}i^{!}\rightarrow \id$ est un
monomorphisme. Les foncteurs $i^{!}$ (pour des anneaux variables)
commutent aux foncteurs restriction des scalaires.
\end{enonce*}

Il est clair que tout morphisme
$$
(0,N;0)\longrightarrow (M,N;N\xrightarrow{u}i^{\ast}j_{\ast}M)
$$
se factorise d'une manière unique par $(0,\ker(u);0)$; ce qui
démontre la propriété d'adjonction. Les autres assertions sont
triviales.

\begin{enonce*}{Proposition 14.6}\etiquette[14.6]{IV.prop14.6}
Pour\pageoriginale tout objet $P$ de ${}_AE$, on a les suites
exactes fonctorielles en $P$:
\begin{gather}
0\longrightarrow j_{!}j^{\ast}P\longrightarrow P\longrightarrow
i_{\ast}i^{\ast}P\longrightarrow 0{\quoi};\tag{14.6.1}\\
0\longrightarrow i_{!}i^{\ast}P\longrightarrow P\longrightarrow
j_{\ast}j^{\ast}P{\quoi},\tag{14.6.2}
\end{gather}
où les flèches non triviales sont les flèches d'adjonction.
\end{enonce*}

Remarquons d'abord qu'une suite
$$
(M', N'; u')\longrightarrow (M,N;u)\longrightarrow (M'',N'';u'')
$$
de $({}_{A/U}U, {}_{A/F}F; i^{\ast}j_{\ast})$ est exacte si et
seulement si les suites correspondantes
\begin{align*}
& M'\longrightarrow M\longrightarrow M''\\
& N'\longrightarrow N\longrightarrow N''
\end{align*}
sont exactes.

Les suites (14.6.1) et (14.6.2) sont transformées par
l'équivalence $\Phi$ en des suites du type \eqhyperref[IV.subsec14.3]{14.3}:
\begin{align*}
& 0\longrightarrow (M,0;0)\longrightarrow (M,N;u)\longrightarrow
(0,N;0)\longrightarrow 0\\
& 0\longrightarrow (0,\ker(u);0)\longrightarrow
(M,N,u)\longrightarrow (M,i^{\ast}j_{\ast}M; \id){\quoi}.
\end{align*}
La vérification de l'exactitude de ces suites est triviale.

\setcounter{subsection}{6}
\subsection{}\etiquette[14.7]{IV.subsec14.7}
La suite exacte (14.6.2) permet d'obtenir une nouvelle
interprétation du foncteur $i_{!}i^{\ast}$. En effet soit $X$ un
objet de $E$. De (14.6.2) on tire la suite exacte de groupes
commutatifs:
$$
0\longrightarrow \Hom_E(X,i_{\ast}i^{!}P)\longrightarrow
\Hom_E(X,P)\longrightarrow \Hom_E(X,j_{\ast}j^{\ast}P){\quoi}.
$$
Notons encore $U$ l'ouvert de $E$, objet final du sous-topos ouvert
$U$. Il résulte des propriétés d'adjonction des foncteurs
$j_{\ast}$, $j^{\ast}$ et $j_{!}^{\ens}$ que le groupe
$\Hom_E(X,j_{\ast}j^{\ast}P)$ est canoniquement isomorphe à
$\Hom_E(X\times U,P)$ et que le morphisme $\Hom_E(X,P)\rightarrow
\Hom_E(X,j_{\ast}j^{\ast}P)$ n'est autre que le morphisme défini par
le monomorphisme $X\times U\rightarrow X$ (\hyperref[IV.sec5]{5}). En
d'autres termes (\hyperref[IV.subsubsec8.5.2]{8.5.2}):

\begin{enonce*}{Proposition 14.8}\etiquette[14.8]{IV.prop14.8}
Les\pageoriginale sections de $i_{\ast}i^{!}P$ sur $E_{/X}$ sont les
sections de $P$ sur $E_{/X}$ dont le support
$(\hyperref[IV.subsubsec9.3.5]{9.3.5})$ contient $F$ \sisi{(notation de $5.9.1$)}{}. En
d'autres termes, $i_{\ast}i^{!}P$ est le plus grand sous-faisceau de
$P$ à support contenu dans $F$.
\end{enonce*}

\setcounter{subsection}{8}
\subsection{}\etiquette[14.9]{IV.subsec14.9}
Par abus de langage, on dit parfois que $i_{\ast}i^{!}P$ est le
sous-Module de $P$ défini par les sections de $P$ à support dans
$F$. Il résulte de \hyperref[IV.rem8.5.3]{8.5.3} que cette terminologie ne fait
qu'étendre aux topos généraux une terminologie
utilisée pour les
topos de faisceaux sur des espaces topologiques.

\begin{enonce*}{Proposition 14.10}\etiquette[14.10]{IV.prop14.10}
\begin{enumerate}
\renewcommand{\labelenumi}{\theenumi)}
\item Le foncteur $P\mapsto i_{\ast}i^{\ast}P$ est isomorphe au
foncteur $P\mapsto i_{\ast}i^{\ast}A\otimes_A P$.

\item Le foncteur $P\mapsto i^{!}i^{\ast}P$ est isomorphe au foncteur
$P\rightarrow \scrHom_A(i_{\ast}i^{\ast}A,P)$.
\end{enumerate}
\end{enonce*}

La suite exacte (14.6.1) s'écrit, dans le cas particulier où
$P$ est le $A$-Module $A$:
\begin{equation*}
0\longrightarrow A_U\longrightarrow A\longrightarrow
i_{\ast}i^{\ast}A\longrightarrow 0\tag{14.10.1}
\end{equation*}\etiquette[14.10.1]{IV.eq14.10.1}
où $A_U$ est le $A$-Module libre engendré par l'ouvert $U$
correspondant au sous-topos ouvert $U$ (9 et 11.3.3). Pour tout
$A$-Module $P$, les $A$-Modules $j_{!}j^{\ast}P$ et
$j_{\ast}j^{\ast}P$ sont canoniquement isomorphes respectivement à
$A_U\otimes_A P$ et $\scrHom_A(A_{U}, P)$ (\hyperref[IV.prop12.3]{12.3}. et
\hyperref[IV.cor12.6]{12.6}). De plus, les morphismes canoniques
$j_{!}j^{\ast}P\rightarrow P$ et $P\rightarrow j_{\ast}j^{\ast}P$
proviennent, modulo ces isomorphismes, du monomorphisme
$A_U\rightarrow A$. On tire donc de la suite exacte \sisi{15.10.1}{\ref{{IV.eq14.10.1}}}, deux
suites exactes (\hyperref[IV.subsec12.2]{12.2} et \hyperref[IV.prop12.11]{12.11}):
\begin{align*}
& j_{!}j^{\ast}P\longrightarrow P\longrightarrow
i_{\ast}i^{\ast}A\otimes_AP\longrightarrow 0{\quoi},\\
& 0\longrightarrow \scrHom_A(i_{\ast}i^{\ast}A, P)\longrightarrow
P\longrightarrow j_{\ast} j^{\ast}P{\quoi};
\end{align*}
d'où les isomorphismes annoncés par comparaison avec (14.6.1)
et (14.6.2).

\begin{thebibliography}{99}\pageoriginale
\bibitem{1} M. Artin et B. Mazur, Homotopy of varieties in the etale
topology, in: Proceedings of a Conference on Local Fields,
Driebergen 1966, Springer.
\bibitem{2} P. Gabriel et G. Zisman, Calculus of fractions and
homotopy theory, Ergebnisse der Mathematik, Bd 35.

\bibitem{3} J. Giraud, Algèbre homologique non commutative:
Grundlehren Springer 1971.

\bibitem{4} A. Grothendieck, sur quelques points d'Algèbre
homologique, Tohoku Math. Journal, (cité [Toh]).

\bibitem{5} A. Grothendieck, Fondements de la Géométrie
Algébrique, (Recueil d'exposés Bourbaki 1957/62), Secrétariat
mathématique, II rue P. Curie Paris.

\bibitem{6} A. Grothendieck, Crystals and the De Rham cohomology
of schemes, (notes by I. Coates et O. Jussila), in: Dix exposés sur
la cohomologie des schémas, North Holland Pub, Cie, 1969.

\bibitem{7} A. Grothendieck, Classes de Chern des
représentations linéaires des groupes discrets, in: Dix exposés sur
la cohomologie des schémas, North Holland Pub. Cie, 1969.

\bibitem{8} R. Godement, Théorie des faisceaux,
Act. Scient. Ind. ${\rm n^{\circ}}$ 1 252, (1958), Hermann (Paris)
(cité [TF]).

\bibitem{9} M. Hakim, Topos annelés et schémas relatifs, Thèse
multigraphiée, Orsay 1967.

\bibitem{10} D. Mumford, Picard groups of moduli problems, in
Arithmetic Algebraic Geometry, Harper's Series in Modern
Mathematics.

\bibitem{11} Nguyen Dinh Ngoc, Thèse Sciences Mathématiques,
Paris 1963, ${\rm n^{\circ}}$ 4 995.

\bibitem{12} J.E. Roos, Distributivité des $\varinjlim$ par
rapports aux $\varprojlim$ des topos

\begin{tabular}{@{}r@{\;}l@{\hspace{1cm}}r@{\;}l}
a) & CR t.259 p. 969-972 & b) & CR T. 259 p. 1605-1608\\
c) & CR t.259 p. 1801-1804 & & (août et septembre 1964).
\end{tabular}

\bibitem{13} P. Deligne - D. Mumford, The irreductibility of
the space of curves of given genus, Pub. math. ${\rm n^{\circ}}$ 36
(1969).

\bibitem{14} B. Mitchell, Theory of categories, Academic Press
(1965).

\bibitem{15} J. Lurie, Higher Topos Theory, Annals of Mathematics Studies, 170. Princeton University Press, 
Princeton, NJ, 2009. xviii+925 pp.

\bibitem{16}  P. Johnstone, Topos theory, 
London Mathematical Society Monographs, Vol. 10. 
Academic Press [Harcourt Brace Jovanovich, Publishers], London-New York, 1977. xxiii+367 pp.

\bibitem{17} F. Borceux, G. Van den Bossche, Structure des topologies
d'un topos, Cahiers de topologie et géométrie différentielle catégoriques,
t. 25, n°1 (1984), p 37-39.

\end{thebibliography}

\ifx\danslelivre\undefined
\end{document}
\fi
